
\documentclass[UTF8,11pt,a4paper]{ctexart}
\usepackage[a4paper,left=23mm,right=23mm,top=24mm,bottom=24mm]{geometry}
\usepackage{fontspec}
\setmainfont{DejaVu Sans}
\setCJKmainfont{Noto Serif CJK SC}
\usepackage{amsmath,amssymb}
\usepackage{xcolor}
\usepackage[colorlinks=true,linkcolor=blue!45!black,unicode=true]{hyperref}
\usepackage{enumitem}
\usepackage{needspace}
\usepackage{fancyhdr}
\setlength{\parindent}{2em}
\setlength{\parskip}{0.25em}
\setlist[enumerate]{leftmargin=2.6em,itemsep=0.4em,topsep=0.5em}
\pagestyle{fancy}
\fancyhf{}
\fancyhead[L]{\small 三千弱水 · 大学化学高阶综合题库}
\fancyhead[R]{\small 待专家审校稿}
\fancyfoot[C]{\thepage}
\setlength{\headheight}{15pt}
\sloppy

\begin{document}
\begin{center}
{\LARGE\bfseries 大学化学高阶综合题 · 第9辑 (83题)}\\[0.8em]
{\large 题目、分步解析与数值核验}\\[0.5em]
{\small \textcolor{red!65!black}{审校稿：未经化学专家核准；禁止直接视为正式试卷或出版物}}
\end{center}
\vspace{1em}
\noindent 本册从仓库结构化 JSONL 自动生成。增补和修改应只编辑源题目数据，并重新执行统一编译。
\par 机械方程残差检查只能说明指定模型的内部数值一致，不保证模型现实可行性。第六辑另须完成跨辑模型去重审查。
\tableofcontents
\clearpage
\section*{第9辑 · 题目及解答}
\addcontentsline{toc}{section}{第9辑}
\subsection*{无机化学}
\addcontentsline{toc}{subsection}{第9辑 / 无机化学}
\Needspace{7\baselineskip}
\subsubsection*{UC-ADV9-I76 · d6八面体晶场配对能与自旋态的反常判别}
\noindent\textbf{题目}\par
八面体d6假想配位中心高自旋构型t2g4eg2（S=2，一对电子），低自旋t2g6（S=0，三对电子）。定义t2g单电子能−0.4Δo、eg为+0.6Δo，每电子对能P=18.50kJ/mol，Δo=22.50kJ/mol。求两态相对球形场重心的总能、低态相对高态差值、高自旋自旋-only磁矩及能级交叉时的Δo。不能忽略配对能数目差异。\par
\medskip\noindent\textbf{分步解答}
\begin{enumerate}
\item t2g和eg的电子数分别为HS(4,2)与LS(6,0)。
\item 高态晶场能4(−0.4Δ)+2(+0.6Δ)=−0.4Δ，低态为6(−0.4Δ)=−2.4Δ。
\item 按一对/三对分别加配对能，EHS=−0.4Δ+P=9.5kJ/mol。
\item ELS=−2.4Δ+3P=1.5kJ/mol。
\item 两态差ELS−EHS=2(P−Δ)=-8kJ/mol，负值表示低自旋更稳定。
\item 交叉要求Δ=P；高自旋S=2时μ=2√[S(S+1)]=4.898979486μB，低自旋按S=0磁矩零。
\item 独立复核：将全部所求量逆代到题设的能量、守恒或跃迁公式；三项无量纲残差：0.00e+0、0.00e+0、0.00e+0。
\end{enumerate}
\noindent\textbf{答案摘要}\par
答案1 = 9.5；答案2 = 1.5；答案3 = -8；答案4 = 4.89897949；答案5 = 18.5\par
\noindent\textbf{模型前提：}
题设假想理想化模型有效，忽略未列相互作用；化学真实性、跨辑原创性及出版资格待审。\par
\noindent\textcolor{red!65!black}{[跨辑去重未完成 · 草稿]} \par
\medskip\hrule\medskip
\Needspace{7\baselineskip}
\subsubsection*{UC-ADV9-I77 · Latimer分段氧化还原与歧化方向的能量法判别}
\noindent\textbf{题目}\par
假想元素X有标准半反应X(III)+e−→X(II)，E1=0.7200V；X(II)+2e−→X(0)，E2=0.3100V。298.15K，忽略非标准态活度影响。求X(III)→X(0)三电子总标准还原电位、对应每mol X的标准自由能，以及3X(II)→2X(III)+X(0)的歧化电势、ΔG与自发性。必须按电子数加和自由能而非直接平均电位。\par
\medskip\noindent\textbf{分步解答}
\begin{enumerate}
\item 第一段自由能−FE1、第二段−2FE2，三电子总和为−F(E1+2E2)。
\item 总还原电位Et=(E1+2E2)/3=0.4466666667V。
\item 总标准自由能ΔG=−3FEt=-129290.345J/mol X。
\item 歧化反应中2X(II)氧化为2X(III)，1个X(II)还原为X(0)，因此电子数2。
\item 歧化电池势Ed=E2−E1=-0.41V，ΔGdis=−2FEd=79117.97234J/(3mol X(II))。
\item Ed<0，故正向歧化在标准态非自发；逆向归中更有利。
\item 独立复核：将全部所求量逆代到题设的能量、守恒或跃迁公式；三项无量纲残差：-1.11e-16、-1.46e-11、0.00e+0。
\end{enumerate}
\noindent\textbf{答案摘要}\par
答案1 = 0.446666667；答案2 = -129290.345；答案3 = -0.41；答案4 = 79117.9723\par
\noindent\textbf{模型前提：}
题设假想理想化模型有效，忽略未列相互作用；化学真实性、跨辑原创性及出版资格待审。\par
\noindent\textcolor{red!65!black}{[跨辑去重未完成 · 草稿]} \par
\medskip\hrule\medskip
\Needspace{7\baselineskip}
\subsubsection*{UC-ADV9-I78 · Pourbaix边界与保护膜电位窗口的参数交点}
\noindent\textbf{题目}\par
假想酸碱耦合电极边界在298.15K为E=E0−0.05916pH（V），E0=1.020V，且一个保护膜允许的最低电位为0.450V。求pH6.500处反应边界电位、保护膜电位与该边界相交的pH、当实际电位0.600V时相对边界的偏移，并指出热力学势线不能直接预测腐蚀电流或钝化动力学。\par
\medskip\noindent\textbf{分步解答}
\begin{enumerate}
\item 电势边界每升高1pH下降0.05916V，这是题给电子/质子计量对应的Nernst线。
\item pH=6.5处边界Eb=1.020−0.05916×6.5=0.63546V。
\item 令Eb=0.450V，交点pH*=(1.020−0.450)/0.05916=9.634888438。
\item 实际电位0.600V高于边界的量为ΔE=0.600−Eb=-0.03546V。
\item 在交点另一侧需要重新判断保护膜窗口与氧化还原稳定区是否交叠。
\item 过电位必须结合电极动力学与表面条件才能对应电流密度，不能仅作平衡电位差外推。
\item 独立复核：将全部所求量逆代到题设的能量、守恒或跃迁公式；三项无量纲残差：0.00e+0、-5.55e-17、0.00e+0。
\end{enumerate}
\noindent\textbf{答案摘要}\par
答案1 = 0.63546；答案2 = 9.63488844；答案3 = -0.03546\par
\noindent\textbf{模型前提：}
题设假想理想化模型有效，忽略未列相互作用；化学真实性、跨辑原创性及出版资格待审。\par
\noindent\textcolor{red!65!black}{[跨辑去重未完成 · 草稿]} \par
\medskip\hrule\medskip
\Needspace{7\baselineskip}
\subsubsection*{UC-ADV9-I79 · 离子晶体的晶胞质量、密度亏损和化学式空位反演}
\noindent\textbf{题目}\par
假想立方晶体AB每式摩尔质量180.0g/mol，一个晶胞有Z=4个AB式单元，晶格参数a=0.5400nm。求晶胞体积、理论晶体密度、0.1000mol AB对应晶胞数；若实测密度比理论密度低2.00\%，晶格常数不变，且只由缺失完整AB式单元解释，求每晶胞缺失式单元平均数及空位比例。说明不排除晶界孔隙或其他缺陷。\par
\medskip\noindent\textbf{分步解答}
\begin{enumerate}
\item 晶格常数由nm化为m：a=5.400×10\textasciicircum{}-10m，晶胞体积V=a³=1.57464e-28m³。
\item 每晶胞质量m=Z MW/NA，式量应使用0.1800kg/mol而非180kg/mol。
\item 理论密度ρ=Z MW/(NA a³)=7592.771226kg/m³。
\item 0.1000mol的AB式单元数为0.1NA，每晶胞Z=4，故晶胞数=1.50553519e+22。
\item 固定a且全由化学式单元空位解释2.00\%质量亏损，每晶胞缺失Z×.02=0.08个式单元。
\item 相对四个AB式位点的缺失比例0.0200；这一结论依赖唯一质量亏损机制假设。
\item 独立复核：将全部所求量逆代到题设的能量、守恒或跃迁公式；三项无量纲残差：0.00e+0、0.00e+0、0.00e+0。
\end{enumerate}
\noindent\textbf{答案摘要}\par
答案1 = 1.57464e-28；答案2 = 7592.77123；答案3 = 1.50553519e+22；答案4 = 0.08；答案5 = 0.02\par
\noindent\textbf{模型前提：}
题设假想理想化模型有效，忽略未列相互作用；化学真实性、跨辑原创性及出版资格待审。\par
\noindent\textcolor{red!65!black}{[跨辑去重未完成 · 草稿]} \par
\medskip\hrule\medskip
\Needspace{7\baselineskip}
\subsubsection*{UC-ADV9-I80 · 光催化半导体带隙吸收与光量子热弛豫预算}
\noindent\textbf{题目}\par
假想无机半导体直接带隙Eg=2.300eV，一束400.0nm单色光入射；h=6.62607015×10\textasciicircum{}-34J·s、c=2.99792458×10\textasciicircum{}8m/s、NA=6.02214076×10\textasciicircum{}23mol\textasciicircum{}-1。假设一个光子激发一个电子且超过带隙能量全部热化。求400nm光子能量(eV与kJ/mol)、热化超额能量(eV)，再判断500nm光子能否单光子直接跨隙。\par
\medskip\noindent\textbf{分步解答}
\begin{enumerate}
\item 用E=hc/λ，400nm转换为4×10\textasciicircum{}-7m，单光子E=4.966114643e-19J。
\item 1eV=1.602176634×10\textasciicircum{}-19J，故光子能量=3.099604961eV。
\item 每摩尔光子能量为NA E/1000=299.0664141kJ/mol。
\item 超过带隙的热化上限E−Eg=0.7996049608eV。
\item 500nm单光子能量=2.479683969eV，低于2.300eV，故题设无缺陷单光子直接跨隙不允许。
\item 若存在杂质态、多光子吸收或声子辅助，则需要改变模型，不能由该简单阈值排除一切吸收。
\item 独立复核：将全部所求量逆代到题设的能量、守恒或跃迁公式；三项无量纲残差：0.00e+0、0.00e+0、0.00e+0。
\end{enumerate}
\noindent\textbf{答案摘要}\par
答案1 = 3.09960496；答案2 = 299.066414；答案3 = 0.799604961；答案4 = 2.47968397\par
\noindent\textbf{模型前提：}
题设假想理想化模型有效，忽略未列相互作用；化学真实性、跨辑原创性及出版资格待审。\par
\noindent\textcolor{red!65!black}{[跨辑去重未完成 · 草稿]} \par
\medskip\hrule\medskip
\Needspace{7\baselineskip}
\subsubsection*{UC-ADV9-I81 · XRD晶面应变对Bragg角与衍射峰位的精确和微分预测}
\noindent\textbf{题目}\par
一阶Bragg XRD测量，Cu Kα波长0.15406nm，未应变时晶面间距d=0.2140nm。晶体出现均匀法向拉伸应变ε=+0.002000，d'=(1+ε)d，波长不变。求原2θ、应变后的2θ、精确峰位移动量以及用一阶微分公式估算的峰位移动，说明2θ与θ差一个因子2且应变无量纲。\par
\medskip\noindent\textbf{分步解答}
\begin{enumerate}
\item 一阶Bragg式2d sinθ=λ，故θ=asin(λ/2d)。
\item 初始2θ=42.19465253度。
\item 应变后d'=d(1+ε)，精确2θ'=42.10642度。
\item 精确移动Δ(2θ)=-0.0882325301度；拉伸增加d导致峰位向低角移动。
\item 微分2d sinθ=λ，得到dθ=−tanθ dd/d；所以Δ(2θ)≈−2tanθε=-0.08842211812度。
\item 微分近似与精确差为0.0001895880209度，ε小使一阶估计合理。
\item 独立复核：将全部所求量逆代到题设的能量、守恒或跃迁公式；三项无量纲残差：0.00e+0、0.00e+0、0.00e+0。
\end{enumerate}
\noindent\textbf{答案摘要}\par
答案1 = 42.1946525；答案2 = 42.10642；答案3 = -0.0882325301；答案4 = -0.0884221181\par
\noindent\textbf{模型前提：}
题设假想理想化模型有效，忽略未列相互作用；化学真实性、跨辑原创性及出版资格待审。\par
\noindent\textcolor{red!65!black}{[跨辑去重未完成 · 草稿]} \par
\medskip\hrule\medskip
\Needspace{7\baselineskip}
\subsubsection*{UC-ADV9-I82 · 轴向g张量EPR谱的定频共振场与微波频率交叉核对}
\noindent\textbf{题目}\par
假想轴向S=1/2配位中心在共振磁场平行分子轴时g\_parallel=2.120、B\_parallel=0.3420T，垂直轴g\_perp=2.010；无零场分裂和超精细分裂。以hν=gμB B，μB=9.2740100783×10\textasciicircum{}-24J/T，求微波频率GHz、垂直共振磁场、场移动量及为什么g更大的方向场更小。\par
\medskip\noindent\textbf{分步解答}
\begin{enumerate}
\item 按简单Zeeman分裂，E=hν=gμB B。
\item 由平行方向计算ν=g\_parallel μB B\_parallel/h=10.14783743GHz。
\item 固定微波频率意味着gparallel Bparallel=gperp Bperp，得Bperp=0.3607164179T。
\item 场移动Bperp−Bparallel=0.01871641791T。
\item 因为gB乘积固定，较小g对应更高的共振磁场。
\item 若存在显著超精细耦合或S>1/2零场分裂，还需完整自旋哈密顿量。
\item 独立复核：将全部所求量逆代到题设的能量、守恒或跃迁公式；三项无量纲残差：0.00e+0、0.00e+0、0.00e+0。
\end{enumerate}
\noindent\textbf{答案摘要}\par
答案1 = 10.1478374；答案2 = 0.360716418；答案3 = 0.0187164179\par
\noindent\textbf{模型前提：}
题设假想理想化模型有效，忽略未列相互作用；化学真实性、跨辑原创性及出版资格待审。\par
\noindent\textcolor{red!65!black}{[跨辑去重未完成 · 草稿]} \par
\medskip\hrule\medskip
\Needspace{7\baselineskip}
\subsubsection*{UC-ADV9-I83 · d8四面体与平方平面构型的能量竞争}
\noindent\textbf{题目}\par
假想d8离子配体两种几何态共享能量零点，四面体晶场总贡献+1.2Δ且有1对电子，平方平面晶场总贡献−0.8Δ且有2对电子。Δ=25.00kJ/mol、每电子对配对能P=18.00kJ/mol。求两个总能、平方平面相对四面体能差和配对能交叉阈值，并解释为何不同几何晶场的参照能必须统一。\par
\medskip\noindent\textbf{分步解答}
\begin{enumerate}
\item 设总能等于晶场能与电子配对能之和，统一能量零点下可比较。
\item 四面体Et=1.2Δ+P=48kJ/mol。
\item 平方平面Es=−0.8Δ+2P=16kJ/mol。
\item Es−Et=P−2Δ=-32kJ/mol，负号表示平方平面更稳定。
\item 能量交叉时P*=2Δ=50kJ/mol。
\item 题给场能未包含配体几何形变、溶剂及配位熵，不能用这条简化能量差推断所有真实d8离子的构型。
\item 独立复算：将所得量重新代回题设的质量作用、能量或位点守恒式；归一化残差：0.00e+0、0.00e+0、0.00e+0。
\end{enumerate}
\noindent\textbf{答案摘要}\par
答案1 = 48；答案2 = 16；答案3 = -32；答案4 = 50\par
\noindent\textbf{模型前提：}
题目采用理想化假想物种，忽略未指明相互作用；尚待人工科学性及模型重复审定。\par
\noindent\textcolor{red!65!black}{[跨辑去重未完成 · 草稿]} \par
\medskip\hrule\medskip
\Needspace{7\baselineskip}
\subsubsection*{UC-ADV9-I84 · Jahn–Teller四方拉伸与八面体键长畸变正交模}
\noindent\textbf{题目}\par
一假想Cu(II)六配位体轴向有2根M–L键2.260Å，赤道4根2.020Å。定义Q3=(2/√6)(2Rax−Req1−Req2)与Q2=(Req1−Req2)/√2，两赤道方向Req1、Req2均为2.020Å。求畸变模式Q3与Q2、六键平均键长、轴向/赤道键长比，并判断几何属于四方还是斜方静态畸变。\par
\medskip\noindent\textbf{分步解答}
\begin{enumerate}
\item 先识别四根赤道键共两对且相等，Req1=Req2=2.020Å。
\item 四方模式Q3=(2/√6)(2×2.260−2.020−2.020)=0.3919183588Å。
\item 斜方模式Q2=(Req1−Req2)/√2=0Å。
\item 六键平均长Rbar=(2Rax+4Req)/6=2.1Å。
\item 轴/赤道键长比Rax/Req=1.118811881；Q3>0而Q2=0，表示静态四方伸长。
\item 光谱若只得到动态平均键长，仍需区分动态Jahn–Teller效应和实际静态构型。
\item 独立复算：将所得量重新代回题设的质量作用、能量或位点守恒式；归一化残差：0.00e+0、1.78e-15、0.00e+0。
\end{enumerate}
\noindent\textbf{答案摘要}\par
答案1 = 0.391918359；答案2 = 0；答案3 = 2.1；答案4 = 1.11881188\par
\noindent\textbf{模型前提：}
题目采用理想化假想物种，忽略未指明相互作用；尚待人工科学性及模型重复审定。\par
\noindent\textcolor{red!65!black}{[跨辑去重未完成 · 草稿]} \par
\medskip\hrule\medskip
\Needspace{7\baselineskip}
\subsubsection*{UC-ADV9-I85 · Born–Haber循环离子形成能与晶格焓的严格符号审计}
\noindent\textbf{题目}\par
假想盐MX(s)形成焓−440.0kJ/mol，对应反应M(s)+0.5X2(g)→MX(s)。给定M升华焓+110.0、X2(g)完全解离2X(g)为+500.0kJ/mol X2、M原子一电离+500.0、X原子电子亲和焓−320.0kJ/mol。求气相离子形成的总焓、气态离子成晶体的晶格生成焓、反向晶格解离焓，并证明为何解离步骤必须取1/2。\par
\medskip\noindent\textbf{分步解答}
\begin{enumerate}
\item 每mol MX只需要1mol X原子，所以X2解离的热量为0.5×500=250kJ/mol。
\item 四项气相步骤总焓Hgas=Hsub+Hatom+Eion+EA=540kJ/mol。
\item Hess闭合Hform=Hgas+Hlattice,form，得到Hlattice,form=-980kJ/mol。
\item 反向拆晶格为气态离子需要Hlattice,diss=980kJ/mol。
\item 回代四项再加晶格生成焓必须恰为实验给出的形成焓−440.0kJ/mol。
\item 若错误把X2解离一整个分子按每mol MX计，结果将差250kJ/mol，属于化学计量错误。
\item 独立复算：将所得量重新代回题设的质量作用、能量或位点守恒式；归一化残差：0.00e+0、0.00e+0、0.00e+0。
\end{enumerate}
\noindent\textbf{答案摘要}\par
答案1 = 540；答案2 = -980；答案3 = 980\par
\noindent\textbf{模型前提：}
题目采用理想化假想物种，忽略未指明相互作用；尚待人工科学性及模型重复审定。\par
\noindent\textcolor{red!65!black}{[跨辑去重未完成 · 草稿]} \par
\medskip\hrule\medskip
\Needspace{7\baselineskip}
\subsubsection*{UC-ADV9-I86 · Marcus外球电子转移的正常区与倒转区界限}
\noindent\textbf{题目}\par
假想无机配合物外球一电子转移在298.15K下重组能λ=1.200eV/电子、标准反应自由能ΔG°=−0.3600eV/电子。采用非绝热Marcus抛物线ΔG‡=(λ+ΔG°)²/(4λ)，且1eV/粒子乘F转为J/mol电子。求摩尔λ与ΔG、活化势垒(kJ/mol)、相对零势垒的Boltzmann因子，以及激活势垒为零的驱动力，并判断继续增加放热性是否永远加速。\par
\medskip\noindent\textbf{分步解答}
\begin{enumerate}
\item λ按1.200F转为摩尔能λmol=115.7823985kJ/mol。
\item 反应自由能ΔGmol=−0.360F/1000=-34.73471956kJ/mol。
\item Marcus势垒ΔG‡=(λ+ΔG)²/(4λ)=14.18334382kJ/mol。
\item 热活化因子exp(−ΔG‡/RT)=0.003274806679。
\item 势垒为零的条件ΔG°=−λ，故临界驱动力−1.200eV/电子。
\item 更强放热的ΔG°低于−λ将进入倒转区，势垒反而升高；真实转移需满足弱耦合及近似谐振反应坐标。
\item 独立复算：将所得量重新代回题设的质量作用、能量或位点守恒式；归一化残差：0.00e+0、0.00e+0、0.00e+0。
\end{enumerate}
\noindent\textbf{答案摘要}\par
答案1 = 115.782399；答案2 = -34.7347196；答案3 = 14.1833438；答案4 = 0.00327480668；答案5 = -1.2\par
\noindent\textbf{模型前提：}
题目采用理想化假想物种，忽略未指明相互作用；尚待人工科学性及模型重复审定。\par
\noindent\textcolor{red!65!black}{[跨辑去重未完成 · 草稿]} \par
\medskip\hrule\medskip
\Needspace{7\baselineskip}
\subsubsection*{UC-ADV9-I87 · EXAFS双壳层散射幅度的Debye–Waller无序修正}
\noindent\textbf{题目}\par
假想金属EXAFS同波数k=12.00Å\textasciicircum{}-1上有两条独立散射路径，第一条占有比例0.650、配位数5.300、σ1=0.00700Å，第二条占比0.350、配位数3.100、σ2=0.01400Å。忽略其他路径因子，定义峰幅Ai=fiNi exp(−2k²σi²)。求两路径峰幅、峰幅比以及若忽略无序指数会错误得到的比值，分析为何单点幅度不能唯一反演配位数。\par
\medskip\noindent\textbf{分步解答}
\begin{enumerate}
\item Debye–Waller衰减是exp(−2k²σ²)，k单位Å\textasciicircum{}-1与σ单位Å使指数无量纲。
\item 第一路径A1=f1N1 exp(−2k²σ1²)=3.396725585。
\item 第二路径A2=f2N2 exp(−2k²σ2²)=1.025450458。
\item 实际峰幅比A1/A2=3.312422906。
\item 忽略热无序会把比值估为f1N1/(f2N2)=3.175115207。
\item 第二路径无序更大，因此在相同波数衰减更强；真实EXAFS还需相位与平均自由程校正。
\item 独立复算：将所得量重新代回题设的质量作用、能量或位点守恒式；归一化残差：0.00e+0、0.00e+0、0.00e+0。
\end{enumerate}
\noindent\textbf{答案摘要}\par
答案1 = 3.39672559；答案2 = 1.02545046；答案3 = 3.31242291；答案4 = 3.17511521\par
\noindent\textbf{模型前提：}
题目采用理想化假想物种，忽略未指明相互作用；尚待人工科学性及模型重复审定。\par
\noindent\textcolor{red!65!black}{[跨辑去重未完成 · 草稿]} \par
\medskip\hrule\medskip
\Needspace{7\baselineskip}
\subsubsection*{UC-ADV9-I88 · Curie–Weiss磁化率中温度无关本底的识别}
\noindent\textbf{题目}\par
假想顺磁晶体体磁化率χ(T)=C/(T−Θ)+χ0，Curie常数C=0.006000K，Weiss温度Θ=270.0K，温度无关本底χ0=2.000×10\textasciicircum{}-5。求300K和340K总χ、扣除本底后的两组顺磁χ、其温度变化比，以及错误按纯Curie律χ=C/T反推300K常数的偏差。\par
\medskip\noindent\textbf{分步解答}
\begin{enumerate}
\item 顺磁部分是C/(T−Θ)而不是C/T。
\item 300K顺磁贡献χpara,300=0.0002，总χ300=0.00022。
\item 340K顺磁贡献χpara,340=0.00008571428571，总χ340=0.0001057142857。
\item 顺磁温度比χpara,340/χpara,300=0.4285714286。
\item 误用Curie律从300K总χ反推Cwrong=300χtotal,300=0.066K，明显偏离给定0.006000K。
\item Θ是平均场描述的交换作用温度尺度，不能直接宣称等于真正有序转变温度。
\item 独立复算：将所得量重新代回题设的质量作用、能量或位点守恒式；归一化残差：0.00e+0、0.00e+0、0.00e+0。
\end{enumerate}
\noindent\textbf{答案摘要}\par
答案1 = 0.00022；答案2 = 0.000105714286；答案3 = 0.0002；答案4 = 8.57142857e-05；答案5 = 0.428571429；答案6 = 0.066\par
\noindent\textbf{模型前提：}
题目采用理想化假想物种，忽略未指明相互作用；尚待人工科学性及模型重复审定。\par
\noindent\textcolor{red!65!black}{[跨辑去重未完成 · 草稿]} \par
\medskip\hrule\medskip
\Needspace{7\baselineskip}
\subsubsection*{UC-ADV9-I89 · Nernst–Einstein固态离子电导与示踪扩散迁移率交叉核算}
\noindent\textbf{题目}\par
假想固态电解质650K单价Li+移动载流子密度8.000×10\textasciicircum{}25m\textasciicircum{}-3，示踪扩散系数D=1.500×10\textasciicircum{}-12m²/s；假设Haven因子为1。基本电荷1.602176634×10\textasciicircum{}-19C，Boltzmann常数1.380649×10\textasciicircum{}-23J/K。求电导率σ、迁移率μ、施加电场150V/m时的电流密度，并由漂移电流j=cqμE独立验证。\par
\medskip\noindent\textbf{分步解答}
\begin{enumerate}
\item 独立载流子Nernst–Einstein式σ=cq²D/(kBT)。
\item 代入数密度、扩散系数、T得σ=0.0003432459283S/m。
\item 迁移率μ=qD/(kBT)=2.67796572e-11m²/(V·s)。
\item 电流密度j=σE=0.05148688925A/m²。
\item 独立采用漂移载流子式j=cqμE，应给出完全相同结果。
\item 真实晶体跳跃常有相关性，若Haven因子不为1，则示踪扩散与电导扩散不一定相等。
\item 独立复算：将所得量重新代回题设的质量作用、能量或位点守恒式；归一化残差：0.00e+0、0.00e+0、0.00e+0。
\end{enumerate}
\noindent\textbf{答案摘要}\par
答案1 = 0.000343245928；答案2 = 2.67796572e-11；答案3 = 0.0514868892\par
\noindent\textbf{模型前提：}
题目采用理想化假想物种，忽略未指明相互作用；尚待人工科学性及模型重复审定。\par
\noindent\textcolor{red!65!black}{[跨辑去重未完成 · 草稿]} \par
\medskip\hrule\medskip
\Needspace{7\baselineskip}
\subsubsection*{UC-ADV9-I90 · McLean晶界偏析等温式的界面富集和有限位点解读}
\noindent\textbf{题目}\par
假想双组元无机晶体B在体相原子分数xb=0.08000，界面理想单占位，B相对体相的偏析自由能ΔGseg=−0.1700eV/原子，温度650K。McLean式xs/(1−xs)=[xb/(1−xb)]exp(−ΔGseg/kBT)。求界面B分数、相对体相富集倍数、晶界B/A比、1000个独立晶界位点的平均B数，并解释若体相有限为何还需总原子守恒。\par
\medskip\noindent\textbf{分步解答}
\begin{enumerate}
\item 偏析因子K=exp(−ΔGseg/kBT)=exp[0.1700F/(RT)]=20.80155651。
\item 由McLean比值求xs=Kxb/[1−xb+Kxb]=0.6439800047。
\item 表面富集倍数xs/xb=8.049750058。
\item B/A位点比xs/(1−xs)=1.808831001，应等于Kxb/(1−xb)。
\item 1000位点的统计期望B个数=643.9800047，不要求为整数。
\item 若晶体体相是有限体积，B的界面富集会耗尽体相，须把体相浓度也作为未知量加入B元素守恒。
\item 独立复算：将所得量重新代回题设的质量作用、能量或位点守恒式；归一化残差：3.55e-15、0.00e+0、0.00e+0。
\end{enumerate}
\noindent\textbf{答案摘要}\par
答案1 = 0.643980005；答案2 = 8.04975006；答案3 = 1.808831；答案4 = 643.980005\par
\noindent\textbf{模型前提：}
题目采用理想化假想物种，忽略未指明相互作用；尚待人工科学性及模型重复审定。\par
\noindent\textcolor{red!65!black}{[跨辑去重未完成 · 草稿]} \par
\medskip\hrule\medskip
\Needspace{7\baselineskip}
\subsubsection*{UC-ADV9-I91 · Mars–van Krevelen固体氧化还原位点的稳态循环产率}
\noindent\textbf{题目}\par
假想氧化物催化表面只含氧化态SO和还原空位SV，SO+SV=1。底物与SO反应还原为SV的条件一级常数kred=0.600s\textasciicircum{}-1，再氧化SV→SO的常数kox=0.200s\textasciicircum{}-1。活性位点总量0.003000mol，产品形成与还原SO一步一一对应。求稳态两种位点比例、总产品速率以及相对无限快速再氧化时的产率比例。\par
\medskip\noindent\textbf{分步解答}
\begin{enumerate}
\item 表面位点连续方程dSO/dt=−kred SO+kox SV，且SO+SV=1。
\item 稳态平衡kred SO=kox SV，故SO=kox/(kred+kox)=0.25。
\item 空位SV=kred/(kred+kox)=0.75。
\item 单位位点周转频率rsite=kred SO=kred kox/(kred+kox)=0.15s\textasciicircum{}-1。
\item 总产品速率R=Csites rsite=0.00045mol/s。
\item 再氧化无限快极限SO→1，速率上限Csites kred，因此当前比例SO=0.25。
\item 独立复算：将所得量重新代回题设的质量作用、能量或位点守恒式；归一化残差：-1.11e-16、0.00e+0、0.00e+0。
\end{enumerate}
\noindent\textbf{答案摘要}\par
答案1 = 0.25；答案2 = 0.75；答案3 = 0.00045；答案4 = 0.25\par
\noindent\textbf{模型前提：}
题目采用理想化假想物种，忽略未指明相互作用；尚待人工科学性及模型重复审定。\par
\noindent\textcolor{red!65!black}{[跨辑去重未完成 · 草稿]} \par
\medskip\hrule\medskip
\Needspace{7\baselineskip}
\subsubsection*{UC-ADV9-I92 · 平方平面Pt(II)配位重排的两条竞争路径}
\noindent\textbf{题目}\par
假想平方平面Pt(II)配合物A在外部过量配体L中分别经两条不可逆并行伪一级反应生成cis与trans产物，速率常数kcis=0.120min\textasciicircum{}-1、ktrans=0.0300min\textasciicircum{}-1，初始A=1.000mmol。求7min未反应A、两异构产物mol数、产物cis比例和瞬时cis生成速率，说明仅靠最终异构体比例并不能证明微观trans effect的机理。\par
\medskip\noindent\textbf{分步解答}
\begin{enumerate}
\item 两个路径共耗底物，故总一级消耗常数k=0.150min\textasciicircum{}-1。
\item 7min残余A=A0exp(−kt)=0.3499377491mmol。
\item 累计cis产物=A0(kcis/k)(1−exp(−kt))=0.5200498007mmol。
\item trans产物=0.1300124502mmol，总产物=1−A。
\item cis选择性C/(C+D)=kcis/k=0.8。
\item 瞬时cis生成速率kcis A(7)=0.04199252989mmol/min，须区别累计产率与瞬时通量。
\item 独立复核：用题设的质量、电荷、波数或者位点比值重算，归一化残差为0.00e+0、-1.11e-16、0.00e+0。
\end{enumerate}
\noindent\textbf{答案摘要}\par
答案1 = 0.349937749；答案2 = 0.520049801；答案3 = 0.13001245；答案4 = 0.8；答案5 = 0.0419925299\par
\noindent\textbf{模型前提：}
所给定理想化学模型用于练习；材料适用性、跨辑原创性与专家核定尚未完成。\par
\noindent\textcolor{red!65!black}{[跨辑去重未完成 · 草稿]} \par
\medskip\hrule\medskip
\Needspace{7\baselineskip}
\subsubsection*{UC-ADV9-I93 · 穆斯堡尔双位点谱的无反冲分数与真实价态占比}
\noindent\textbf{题目}\par
某假想混合价矿物的穆斯堡尔谱Fe(III)、Fe(II)两个峰簇未归一化积分面积分别0.5100与0.1920，对应无反冲因子f3=0.8500、f2=0.4800，薄层线性吸收近似成立。求两铁环境真实原子比例、平均铁价、未经无反冲校正的表观Fe(III)比例，并解释峰面积与原子数不成比例的原因。\par
\medskip\noindent\textbf{分步解答}
\begin{enumerate}
\item 薄层近似谱面积Ai与该环境铁原子数Ni及无反冲分数fi之积成正比。
\item 校正权重w3=0.510/0.850=0.6。
\item 校正权重w2=0.192/0.480=0.4。
\item 真实Fe(III)比例x=w3/(w3+w2)=0.6，Fe(II)比例=0.4。
\item 平均Fe价=3x+2(1−x)=2.6。
\item 若直接按谱面积归一，则错误Fe(III)比例=0.7264957265，偏差来自两环境不同无反冲因子。
\item 独立复核：用题设的质量、电荷、波数或者位点比值重算，归一化残差为0.00e+0、0.00e+0、0.00e+0。
\end{enumerate}
\noindent\textbf{答案摘要}\par
答案1 = 0.6；答案2 = 0.4；答案3 = 2.6；答案4 = 0.726495726\par
\noindent\textbf{模型前提：}
所给定理想化学模型用于练习；材料适用性、跨辑原创性与专家核定尚未完成。\par
\noindent\textcolor{red!65!black}{[跨辑去重未完成 · 草稿]} \par
\medskip\hrule\medskip
\Needspace{7\baselineskip}
\subsubsection*{UC-ADV9-I94 · Raman反Stokes与Stokes声子热布居温度计}
\noindent\textbf{题目}\par
假想无机晶格Raman峰波数560.0cm\textasciicircum{}-1，在300K下模型Ianti/IStokes=exp(−hcνtilde/kBT)，忽略入射与散射频率四次方修正。h=6.62607015×10\textasciicircum{}-34J·s、c=2.99792458×10\textasciicircum{}8m/s、kB=1.380649×10\textasciicircum{}-23J/K。求声子能量、能量与kBT比、Bose平均声子数及峰强比，并给出由实验峰强比反算温度的公式。\par
\medskip\noindent\textbf{分步解答}
\begin{enumerate}
\item 把560cm\textasciicircum{}-1变成56000m\textasciicircum{}-1，光学振动频率ν=c×56000=16788377650000Hz。
\item 能量E=hν=1.11240968e-20J。
\item 无量纲E/(kBT)=2.685716838。
\item Bose平均占据n=1/(exp(x)−1)=0.07315977896。
\item 反Stokes与Stokes强度比r=n/(n+1)=exp(−x)=0.06817230798。
\item 反演温度T=E/[kB ln(1/r)]；真实Raman测温还需校正两谱线响应及频率因子。
\item 独立复核：用题设的质量、电荷、波数或者位点比值重算，归一化残差为1.39e-17、0.00e+0、0.00e+0。
\end{enumerate}
\noindent\textbf{答案摘要}\par
答案1 = 1.11240968e-20；答案2 = 2.68571684；答案3 = 0.073159779；答案4 = 0.068172308\par
\noindent\textbf{模型前提：}
所给定理想化学模型用于练习；材料适用性、跨辑原创性与专家核定尚未完成。\par
\noindent\textcolor{red!65!black}{[跨辑去重未完成 · 草稿]} \par
\medskip\hrule\medskip
\Needspace{7\baselineskip}
\subsubsection*{UC-ADV9-I95 · 钙钛矿A/B离子半径与双重几何容差区间}
\noindent\textbf{题目}\par
假想ABO3钙钛矿三种离子半径rA=1.400Å、rB=0.620Å、rO=1.350Å。定义Goldschmidt容差因子t=(rA+rO)/(√2(rB+rO))、八面体比μ=rB/rO。求t、μ、t−1并依据题给经验区间0.80<t<1.05且0.414<μ<0.732判断是否满足几何条件，解释几何稳定区间只是经验，不代表必然热力学稳定。\par
\medskip\noindent\textbf{分步解答}
\begin{enumerate}
\item A–O半径和rA+rO=2.75Å。
\item B–O半径和rB+rO=1.97Å。
\item 容差因子t=0.987077994。
\item 八面体半径比μ=0.4592592593。
\item 与理想t=1的偏差t−1=-0.01292200596，逐一检查题定两项几何区间。
\item 真实相稳定还依赖氧价态、晶格自由能和温压条件，不能由简单几何因子作存在性定理。
\item 独立复核：用题设的质量、电荷、波数或者位点比值重算，归一化残差为0.00e+0、0.00e+0、0.00e+0。
\end{enumerate}
\noindent\textbf{答案摘要}\par
答案1 = 0.987077994；答案2 = 0.459259259；答案3 = -0.012922006\par
\noindent\textbf{模型前提：}
所给定理想化学模型用于练习；材料适用性、跨辑原创性与专家核定尚未完成。\par
\noindent\textcolor{red!65!black}{[跨辑去重未完成 · 草稿]} \par
\medskip\hrule\medskip
\Needspace{7\baselineskip}
\subsubsection*{UC-ADV9-I96 · 多核羟桥配离子的核电荷及成桥释放质子当量}
\noindent\textbf{题目}\par
假想配离子[M2(μ-OH)3L2]中金属均为+III，三根桥羟基各−I，两个L配体各−I，且按题设每形成一根羟桥释放一个H+。现有配合物0.01000mol、需要用50.00mL强碱全部中和形成桥时释放的质子。求裸配离子电荷、每mol配离子释放H+当量、所释放H+物质的量及强碱所需浓度，判断把裸配离子写成电中性是否漏反离子。\par
\medskip\noindent\textbf{分步解答}
\begin{enumerate}
\item 两金属的总电荷2×(+3)=+6。
\item 三羟桥总电荷−3、两个L总电荷−2，所以配离子电荷+1。
\item 每mol配离子形成三根羟桥，释放质子3mol H+。
\item 0.01000mol配离子形成时释放nH=3×0.01000=0.03mol。
\item 以0.05000L强碱中和，cOH=nH/V=0.6M。
\item 若把此裸配离子标作电中性，必缺少一个一价反离子或额外负电配体。
\item 独立复核：用题设的质量、电荷、波数或者位点比值重算，归一化残差为0.00e+0、0.00e+0、0.00e+0。
\end{enumerate}
\noindent\textbf{答案摘要}\par
答案1 = 1；答案2 = 3；答案3 = 0.03；答案4 = 0.6\par
\noindent\textbf{模型前提：}
所给定理想化学模型用于练习；材料适用性、跨辑原创性与专家核定尚未完成。\par
\noindent\textcolor{red!65!black}{[跨辑去重未完成 · 草稿]} \par
\medskip\hrule\medskip
\Needspace{7\baselineskip}
\subsubsection*{UC-ADV9-I97 · 双钙钛矿B/B'长程有序衍射强度的平方根反演}
\noindent\textbf{题目}\par
假想A2BB'O6双钙钛矿的B/B'长程有序参数S∈[0,1]，归一化超结构XRD峰强度按Isuper=Iref S²变化；参照Iref=0.04000、实测I=0.01960，吸收和热因子已校正。若各亚晶格反占位概率x=(1−S)/2，求有序参数、反占位率、1000个B位点中的平均错占个数和误用线性强度比估计S的误差。\par
\medskip\noindent\textbf{分步解答}
\begin{enumerate}
\item 结构因子振幅与有序参数S成正比，峰强度为振幅平方。
\item 因此S=√(I/Iref)=0.7。
\item 反占位概率x=(1−S)/2=0.15。
\item 1000个位点的平均错占B个数1000x=150。
\item 若错用Slinear=I/Iref=0.49，则S低估量为0.21。
\item 真实有序参数仍需排除峰重叠、晶粒择优取向和各向异性热振动。
\item 逆算题设谱学、动力学或物料条件的残差：0.00e+0、0.00e+0、0.00e+0。
\end{enumerate}
\noindent\textbf{答案摘要}\par
答案1 = 0.7；答案2 = 0.15；答案3 = 150；答案4 = 0.21\par
\noindent\textbf{模型前提：}
仅题设理想化模型成立，尚未经专家科学审稿或模型原创性核验。\par
\noindent\textcolor{red!65!black}{[跨辑去重未完成 · 草稿]} \par
\medskip\hrule\medskip
\Needspace{7\baselineskip}
\subsubsection*{UC-ADV9-I98 · 快平衡配位物的同位素标签交换与微观速率区分}
\noindent\textbf{题目}\par
假想含同位素标签的无机络合物可快速互变X⇌XL，给定配位平衡XL/X=0.7000且自由配体由外库固定；X与XL分别以kX=0.220h\textasciicircum{}-1和kXL=0.05000h\textasciicircum{}-1不可逆丢失标签，失标签物不计入信号。初始标签量1mol。求有效一级衰减常数、6h剩余标签mol、自由态丢失标签通量贡献及总丢失量，解释为何只测总衰减不能唯一反演两个微观速率。\par
\medskip\noindent\textbf{分步解答}
\begin{enumerate}
\item 快平衡时自由态份额pX=1/(1+0.7)=0.5882352941。
\item 结合态份额pXL=0.7/(1+0.7)=0.4117647059。
\item 观察有效一级常数kobs=pXkX+pXLkXL=0.15h\textasciicircum{}-1。
\item 6h剩余标签L=e\textasciicircum{}(−kobs t)=0.4065696597mol。
\item 自由态贡献份额pXkX/kobs=0.862745098，总标签损失1−L=0.5934303403mol。
\item 若交换并非远快于标签丢失，需解多态耦合微分方程，单一加权常数不成立。
\item 逆算题设谱学、动力学或物料条件的残差：1.04e-17、0.00e+0、0.00e+0。
\end{enumerate}
\noindent\textbf{答案摘要}\par
答案1 = 0.15；答案2 = 0.40656966；答案3 = 0.862745098；答案4 = 0.59343034\par
\noindent\textbf{模型前提：}
仅题设理想化模型成立，尚未经专家科学审稿或模型原创性核验。\par
\noindent\textcolor{red!65!black}{[跨辑去重未完成 · 草稿]} \par
\medskip\hrule\medskip
\Needspace{7\baselineskip}
\subsubsection*{UC-ADV9-I99 · 多晶面无机催化剂的位点分布与TOF贡献差异}
\noindent\textbf{题目}\par
假想两类催化表面(100)和(111)的总可及位点比例分别80.00\%和20.00\%，每位点周转频率分别2.500s\textasciicircum{}-1和0.5000s\textasciicircum{}-1，整体位点数量0.01000mol，稳态下没有交叉转化。求整体平均TOF、总产品生成mol/s、(100)晶面产率分数及其对几何位点占比的富集倍数，并说明位点多不必等于产物贡献按同一比例。\par
\medskip\noindent\textbf{分步解答}
\begin{enumerate}
\item 两类可及位点分数之和为1，均按相同活性单位计。
\item 每位点平均TOF=θ100 k100+θ111 k111=2.1s\textasciicircum{}-1。
\item 整体产品生成速率r=ntotal TOF=0.021mol/s。
\item (100)对总速率贡献分数f100=(θ100k100)/TOF=0.9523809524。
\item 相对于其位点占比0.8的反应贡献富集倍数f100/θ100=1.19047619。
\item 若表面活性位点会相互转化，简单静态加权失效，应增加位点转移动力学。
\item 逆算题设谱学、动力学或物料条件的残差：0.00e+0、0.00e+0、0.00e+0。
\end{enumerate}
\noindent\textbf{答案摘要}\par
答案1 = 2.1；答案2 = 0.021；答案3 = 0.952380952；答案4 = 1.19047619\par
\noindent\textbf{模型前提：}
仅题设理想化模型成立，尚未经专家科学审稿或模型原创性核验。\par
\noindent\textcolor{red!65!black}{[跨辑去重未完成 · 草稿]} \par
\medskip\hrule\medskip
\Needspace{7\baselineskip}
\subsubsection*{UC-ADV9-I100 · M2L3无机自组装体的双元素限量与转换率计量}
\noindent\textbf{题目}\par
假想二金属三配体笼由2M+3L→M2L3形成，初始金属M总量0.08000mol，单齿配体L总量0.1200mol，按金属计的转化率55.00\%，总体积保持0.02000L，无其它物种。求理论最多可生成笼mol、实际生成笼mol、剩余M/L的mol以及笼浓度，解释不可把单一配体的“55\%转化”错误等同于笼的初始金属数。\par
\medskip\noindent\textbf{分步解答}
\begin{enumerate}
\item 按组装计量每个笼消耗2M、3L，理论最大nmax=min(0.080/2,0.120/3)=0.04mol。
\item 金属转化55\%，意味着实际笼形成n=0.55nmax=0.022mol。
\item 消耗M为2n=0.044mol，剩余M=0.036mol。
\item 消耗L为3n=0.066mol，剩余L=0.054mol。
\item 笼浓度为n/0.02000=1.1mol/L，金属及配体两个元素守恒必须同时核对。
\item 如果存在单金属中间体或寡聚副产物，必须建立更完整的物料平衡与配位平衡。
\item 逆算题设谱学、动力学或物料条件的残差：0.00e+0、0.00e+0、0.00e+0。
\end{enumerate}
\noindent\textbf{答案摘要}\par
答案1 = 0.04；答案2 = 0.022；答案3 = 0.036；答案4 = 0.054；答案5 = 1.1\par
\noindent\textbf{模型前提：}
仅题设理想化模型成立，尚未经专家科学审稿或模型原创性核验。\par
\noindent\textcolor{red!65!black}{[跨辑去重未完成 · 草稿]} \par
\medskip\hrule\medskip
\subsection*{分析化学}
\addcontentsline{toc}{subsection}{第9辑 / 分析化学}
\Needspace{7\baselineskip}
\subsubsection*{UC-ADV9-A76 · 毛细管电泳EOF与目标离子有效电泳迁移率的分离}
\noindent\textbf{题目}\par
一支有效总长60.00cm、检测器距进样口50.00cm的毛细管电泳管，在20.00kV电压下中性EOF标记峰迁移时间4.000min，阳离子目标峰3.000min。毛细管全长决定电场E=V/L，检测距离决定迁移速度v=Ldet/t。求电渗迁移率μEOF、目标总迁移率μtot、离子相对电泳迁移率μep及若将有效长度误当总长所造成的迁移率相对偏差。\par
\medskip\noindent\textbf{分步解答}
\begin{enumerate}
\item 电场E=20000/0.6000=33333.33333V/m。
\item 由v=Ldet/t，迁移率μ=LdetLtot/(V t)。
\item 中性标记仅随EOF移动，μEOF=6.25e-8m²/(V s)。
\item 阳离子总迁移率μtot=8.333333333e-8m²/(V s)。
\item 分离电泳迁移率μep=μtot−μEOF=2.083333333e-8m²/(V s)。
\item 若错设Ldet=Ltot，会把所有迁移率乘Ltot/Ldet=1.2，高估20\%。
\item 独立将数据代回迁移、保留、能量或响应公式，检查量纲与可逆解的数值残差：0.00e+0、0.00e+0、0.00e+0。
\end{enumerate}
\noindent\textbf{答案摘要}\par
答案1 = 6.25e-08；答案2 = 8.33333333e-08；答案3 = 2.08333333e-08；答案4 = 1.2\par
\noindent\textbf{模型前提：}
题设校准和理想分析模型有效，尚未通过真实仪器实验、跨辑模型原创或专家质量审核。\par
\noindent\textcolor{red!65!black}{[跨辑去重未完成 · 草稿]} \par
\medskip\hrule\medskip
\Needspace{7\baselineskip}
\subsubsection*{UC-ADV9-A77 · 程序升温以外等温GC保留指数的死时间修正}
\noindent\textbf{题目}\par
等温气相色谱中n-癸烷(C10)保留时间11.00min、n-十一烷(C11)17.00min、未知分析物14.00min，未保留标记死时间1.000min。等温Kováts保留指数I=1000+100ln(tR'X/tR'C10)/ln(tR'C11/tR'C10)，其中tR'=tR−tM。求未知分析物指数、三个校正保留时间、未校正用原始tR所得指数与误差，并指出此关系在温度程序GC需要另行定义保留指数。\par
\medskip\noindent\textbf{分步解答}
\begin{enumerate}
\item 保留指数需要用校正保留时间tR'=tR−tM。
\item C10、C11的tR'分别为10、16min，分析物为13min。
\item 代入等温指数I=1055.821753。
\item 若误直接代原始时间，得到Iwrong=1055.399046。
\item 二者差为-0.4227072813，死时间不是可无视的常量平移。
\item 若未知物保留不在两个烷烃之间，应外推并警惕固定相选择性差异。
\item 独立将数据代回迁移、保留、能量或响应公式，检查量纲与可逆解的数值残差：1.11e-16、0.00e+0、0.00e+0。
\end{enumerate}
\noindent\textbf{答案摘要}\par
答案1 = 1055.82175；答案2 = 10；答案3 = 16；答案4 = 13；答案5 = 1055.39905\par
\noindent\textbf{模型前提：}
题设校准和理想分析模型有效，尚未通过真实仪器实验、跨辑模型原创或专家质量审核。\par
\noindent\textcolor{red!65!black}{[跨辑去重未完成 · 草稿]} \par
\medskip\hrule\medskip
\Needspace{7\baselineskip}
\subsubsection*{UC-ADV9-A78 · LC两组分选择性与柱效对分离度的非线性贡献}
\noindent\textbf{题目}\par
等度液相色谱中两个近邻峰保留因子k1=4.000、k2=5.000，理论塔板数N=10000，死时间t0=1.200min。采用近似分离度公式Rs=(√N/4)[(α−1)/α][k2/(1+k2)]，α=k2/k1。求两峰保留时间、选择性α、实际Rs和若要求Rs≥1.500的最小理论塔板数；说明提高k2并不能线性提高Rs。\par
\medskip\noindent\textbf{分步解答}
\begin{enumerate}
\item tR,i=t0(1+ki)，得到tR1=6min、tR2=7.2min。
\item 选择性α=k2/k1=1.25。
\item 将α与k2带入分离度公式得Rs=4.166666667。
\item 令目标Rs=1.500，反解Nmin=[4Rs/coef]²=1296。
\item 塔板数必须为正，严格达到阈值时上取整数=1296。
\item 保留项k2/(1+k2)在高k时趋于1，所以更高保留不保证同比例提高分离度。
\item 独立将数据代回迁移、保留、能量或响应公式，检查量纲与可逆解的数值残差：0.00e+0、0.00e+0、0.00e+0。
\end{enumerate}
\noindent\textbf{答案摘要}\par
答案1 = 6；答案2 = 7.2；答案3 = 1.25；答案4 = 4.16666667；答案5 = 1296\par
\noindent\textbf{模型前提：}
题设校准和理想分析模型有效，尚未通过真实仪器实验、跨辑模型原创或专家质量审核。\par
\noindent\textcolor{red!65!black}{[跨辑去重未完成 · 草稿]} \par
\medskip\hrule\medskip
\Needspace{7\baselineskip}
\subsubsection*{UC-ADV9-A79 · 离子迁移谱Mason–Schamp迁移率与碰撞截面反演}
\noindent\textbf{题目}\par
离子迁移谱IMS中假想单价离子质量200amu、载气N2质量28amu，T=300.0K、压强101325Pa，定态离子-中性碰撞截面Ω=2.000×10\textasciicircum{}-18m²。假定低场Mason–Schamp公式K=(3q/(16NΩ))sqrt[2π/(μkBT)]，其中N=P/(kBT)、μ为约化质量，1amu=1.66053906660×10\textasciicircum{}-27kg。求约化质量、气体分子数密度、迁移率及电场1000V/m下离子漂移速度。\par
\medskip\noindent\textbf{分步解答}
\begin{enumerate}
\item 约化质量μ=mion mgas/(mion+mgas)=4.078517006e-26kg。
\item 理想气体分子数密度N=P/(kBT)=2.446313292e+25m\textasciicircum{}-3。
\item 代入Mason–Schamp式得迁移率K=0.0001184148371m²/(V s)。
\item 电场1000V/m产生漂移速度v=KE=0.1184148371m/s。
\item 要求低场离子漂移及有效截面独立于场强，实际高场IMS不可直接套用。
\item 注意N是分子数密度而非mol/m³，否则量纲会错误。
\item 独立将数据代回迁移、保留、能量或响应公式，检查量纲与可逆解的数值残差：0.00e+0、0.00e+0、0.00e+0。
\end{enumerate}
\noindent\textbf{答案摘要}\par
答案1 = 4.07851701e-26；答案2 = 2.44631329e+25；答案3 = 0.000118414837；答案4 = 0.118414837\par
\noindent\textbf{模型前提：}
题设校准和理想分析模型有效，尚未通过真实仪器实验、跨辑模型原创或专家质量审核。\par
\noindent\textcolor{red!65!black}{[跨辑去重未完成 · 草稿]} \par
\medskip\hrule\medskip
\Needspace{7\baselineskip}
\subsubsection*{UC-ADV9-A80 · ICP-MS双同位素的质量歧视指数校准与未知样校正}
\noindent\textbf{题目}\par
ICP-MS用标准物质的同位素比Rtrue=I62/I60=0.4000校准，仪器测得Rmeas=0.4030；设质量歧视指数律Rmeas=Rtrue(62/60)\textasciicircum{}β。另一未知样的同位素比实测0.5200，且质量歧视指数相同。求β、未知样校正后真实同位素比、质量数62同位素分数与未校正的表观分数，检验误差方向。\par
\medskip\noindent\textbf{分步解答}
\begin{enumerate}
\item 标准样质量歧视因子Fmass=0.4030/0.4000=1.0075。
\item 由Rmeas/Rtrue=(62/60)\textasciicircum{}β取对数求β=0.2278760358。
\item 未知样真实比Rsample=0.5200/Fmass=0.5161290323。
\item 只有两种同位素时62的分数x62=Rsample/(1+Rsample)=0.3404255319。
\item 不校正的表观分数0.52/1.52=0.3421052632，高估真实值。
\item 如果仪器质量歧视非固定指数或样品含同量异位素干扰，需要附加校准信息。
\item 独立将数据代回迁移、保留、能量或响应公式，检查量纲与可逆解的数值残差：0.00e+0、0.00e+0、0.00e+0。
\end{enumerate}
\noindent\textbf{答案摘要}\par
答案1 = 0.227876036；答案2 = 0.516129032；答案3 = 0.340425532；答案4 = 0.342105263\par
\noindent\textbf{模型前提：}
题设校准和理想分析模型有效，尚未通过真实仪器实验、跨辑模型原创或专家质量审核。\par
\noindent\textcolor{red!65!black}{[跨辑去重未完成 · 草稿]} \par
\medskip\hrule\medskip
\Needspace{7\baselineskip}
\subsubsection*{UC-ADV9-A81 · FTIR同位素位移与谐振近似振动零点能差}
\noindent\textbf{题目}\par
一假想CO键近似为同一谐力常数的双原子振子，12C16O基频2140cm\textasciicircum{}-1，同位素替换为13C16O后只改变约化质量。原子质量取整数amu（12、13、16），振动频率与约化质量负平方根成正比。求两种同位素约化质量、替代后的基频、两基频差及1mol CO的零点能变化kJ/mol；假设谐振且力常数保持不变。\par
\medskip\noindent\textbf{分步解答}
\begin{enumerate}
\item 12C16O约化质量μ1=12×16/(12+16)=6.857142857amu。
\item 13C16O约化质量μ2=13×16/(13+16)=7.172413793amu。
\item 同力常数谐振给ν2=ν1sqrt(μ1/μ2)=2092.438506cm\textasciicircum{}-1。
\item 频移ν1−ν2=47.56149407cm\textasciicircum{}-1；重同位素降低振动频率。
\item 每分子零点能差1/2hc(ν1−ν2)，再乘阿伏伽德罗常数得差=0.2844809096kJ/mol。
\item 真实振动存在非谐效应和电子结构的同位素微小改变，这里按题设忽略。
\item 独立将数据代回迁移、保留、能量或响应公式，检查量纲与可逆解的数值残差：0.00e+0、0.00e+0、0.00e+0。
\end{enumerate}
\noindent\textbf{答案摘要}\par
答案1 = 6.85714286；答案2 = 7.17241379；答案3 = 2092.43851；答案4 = 47.5614941；答案5 = 0.28448091\par
\noindent\textbf{模型前提：}
题设校准和理想分析模型有效，尚未通过真实仪器实验、跨辑模型原创或专家质量审核。\par
\noindent\textcolor{red!65!black}{[跨辑去重未完成 · 草稿]} \par
\medskip\hrule\medskip
\Needspace{7\baselineskip}
\subsubsection*{UC-ADV9-A82 · FRET寿命猝灭与分子间距第六次方反演}
\noindent\textbf{题目}\par
荧光标记供体在无受体时寿命τD=4.000ns，受体共存时寿命τDA=1.500ns。假设无其它淬灭、Förster半径R0=5.000nm、距离分布可视为单一固定距离。FRET效率E=1−τDA/τD=1/[1+(r/R0)\textasciicircum{}6]。求转移效率、供受体距离、相对强度损失和如果漏掉六次方而用线性距离假设得到的错误距离。\par
\medskip\noindent\textbf{分步解答}
\begin{enumerate}
\item 由寿命比τDA/τD=1.5/4=0.375得FRET效率E=1−0.375=0.625。
\item Förster关系变形为(r/R0)\textasciicircum{}6=(1/E)−1。
\item 所以r=R0[(1/E)−1]\textasciicircum{}(1/6)=4.591929511nm。
\item 供体寿命与强度在辐射速率不变时比值为τDA/τD=0.375。
\item 错误用线性r\_wrong=R0[(1/E)−1]=3nm，与六次方非线性关系差异明显。
\item 若供受体距离有分布或存在静态淬灭，简单单距离寿命反演会产生偏差。
\item 独立将数据代回迁移、保留、能量或响应公式，检查量纲与可逆解的数值残差：0.00e+0、0.00e+0、0.00e+0。
\end{enumerate}
\noindent\textbf{答案摘要}\par
答案1 = 0.625；答案2 = 4.59192951；答案3 = 0.375；答案4 = 3\par
\noindent\textbf{模型前提：}
题设校准和理想分析模型有效，尚未通过真实仪器实验、跨辑模型原创或专家质量审核。\par
\noindent\textcolor{red!65!black}{[跨辑去重未完成 · 草稿]} \par
\medskip\hrule\medskip
\Needspace{7\baselineskip}
\subsubsection*{UC-ADV9-A83 · ESI-MS多电荷态包络与同位素峰距的中性质量反演}
\noindent\textbf{题目}\par
电喷雾质谱给某蛋白质片段两个相邻电荷态峰簇，质量数与电荷比观测分别为m/z约601.007276和451.007276，两个峰簇内同位素峰间距分别约0.3333和0.2500Th。假设离子形式为[M+zH]\textasciicircum{}\{z+\}，质子质量1.007276Da。求两个电荷数z、蛋白中性质量M和两峰m/z差，说明不能直接把m/z读数作为中性分子量。\par
\medskip\noindent\textbf{分步解答}
\begin{enumerate}
\item 同位素峰间距约为1/z Th，因此0.3333与0.2500分别对应z1=3、z2=4。
\item 带z个质子时m/z=(M+zH)/z=M/z+H。
\item 两峰差Δ=(M/3+H)−(M/4+H)=M/12=150Th。
\item 反解中性质量M=Δ/(1/3−1/4)=1800Da。
\item 将M=1800Da及H=1.007276Da代回两峰m/z，分别得601.007276与451.007276。
\item 若峰来自钠加合、溶剂加合或不同碎片，仅凭这组峰无法按相同H加成式恢复M。
\item 独立将数据代回迁移、保留、能量或响应公式，检查量纲与可逆解的数值残差：0.00e+0、0.00e+0、1.14e-12。
\end{enumerate}
\noindent\textbf{答案摘要}\par
答案1 = 3；答案2 = 4；答案3 = 1800；答案4 = 150；答案5 = 601.007276；答案6 = 451.007276\par
\noindent\textbf{模型前提：}
题设校准和理想分析模型有效，尚未通过真实仪器实验、跨辑模型原创或专家质量审核。\par
\noindent\textcolor{red!65!black}{[跨辑去重未完成 · 草稿]} \par
\medskip\hrule\medskip
\Needspace{7\baselineskip}
\subsubsection*{UC-ADV9-A84 · 计时安培Cottrell扩散电流与恒定背景同时分离}
\noindent\textbf{题目}\par
平面电极有效面积0.1000cm²（即1.000×10\textasciicircum{}-5m²），一电子可逆反应物初始浓度2.000mmol/L，扩散系数D=1.000×10\textasciicircum{}-9m²/s，F=96485.33212C/mol。假设计时安培响应i(t)=nFAC√[D/(πt)]+ib，其中ib是未知恒定背景。实验在t=1.000s和4.000s的电流为由真实背景ib=80.00μA生成的数据。求两个实测电流、由两点逆反演的背景和Cottrell系数，以及忽略背景导致的1秒浓度高估倍数。\par
\medskip\noindent\textbf{分步解答}
\begin{enumerate}
\item 注意浓度2.000mmol/L=2.000mol/m³、面积0.1000cm²=1.000×10\textasciicircum{}-5m²。
\item 定义Cottrell系数a=nFAC√(D/π)=0.0000344283616A·s\textasciicircum{}1/2。
\item t=1s电流i1=a+ib=0.000114428362A，t=4s电流i4=a/2+ib=0.0000972141808A。
\item 两次读数相减给a=2(i1−i4)=0.0000344283616，背景ib=2i4−i1=0.00008A。
\item 若不扣背景，1秒按i1=a假设估浓度会高估i1/a=3.32366562倍。
\item Cottrell仅适用于半无限扩散控制、理想阶跃和早期双电层可忽略或独立估计的区间。
\item 独立用校准式、通量方程或概率守恒重算全部关键结论，归一化残差：-1.36e-20、0.00e+0、-2.03e-20。
\end{enumerate}
\noindent\textbf{答案摘要}\par
答案1 = 0.000114428362；答案2 = 9.72141808e-05；答案3 = 8e-05；答案4 = 3.44283616e-05；答案5 = 3.32366562\par
\noindent\textbf{模型前提：}
理想分析模型与题给标准态成立；不代表完成实际计量溯源、版权或专家科学审校。\par
\noindent\textcolor{red!65!black}{[跨辑去重未完成 · 草稿]} \par
\medskip\hrule\medskip
\Needspace{7\baselineskip}
\subsubsection*{UC-ADV9-A85 · Randles–Ševčík可逆电极峰电流的扫描速率平方根核验}
\noindent\textbf{题目}\par
25℃理想可逆一电子电极反应，用Randles–Ševčík公式ip(A)=2.69×10\textasciicircum{}5 n\textasciicircum{}(3/2) Aelec(cm²) D\textasciicircum{}(1/2)(cm/s)\textasciicircum{}1/2 C(mol/cm³) v\textasciicircum{}(1/2)(V/s)\textasciicircum{}1/2。给定n=1、电极面积0.0700cm²、D=1.000×10\textasciicircum{}-5cm²/s、C=1.000mmol/L、初扫描速率0.1000V/s。求峰电流、扫描速率升至0.4000V/s的峰电流、两者比值与若错误将mM直接当mol/cm³导致的倍数误差。\par
\medskip\noindent\textbf{分步解答}
\begin{enumerate}
\item 先换算浓度1.000mmol/L=1.000×10\textasciicircum{}-6mol/cm³。
\item 把n=1、A=.07、D=10\textasciicircum{}-5、C=10\textasciicircum{}-6及v=.1代入得ip1=0.00001883A。
\item v2=.4000V/s，电流ip2=coef√v2=0.00003766A。
\item 两次电流比ip2/ip1=√(v2/v1)=2。
\item 若把1mM错看作1mol/cm³，将峰电流放大10\textasciicircum{}6倍，属于严重单位错误。
\item 此标准式依赖25℃可逆扩散控制，不可套到强吸附或严重不可逆电子转移动力学。
\item 独立用校准式、通量方程或概率守恒重算全部关键结论，归一化残差：0.00e+0、0.00e+0、0.00e+0。
\end{enumerate}
\noindent\textbf{答案摘要}\par
答案1 = 1.883e-05；答案2 = 3.766e-05；答案3 = 2；答案4 = 1000000\par
\noindent\textbf{模型前提：}
理想分析模型与题给标准态成立；不代表完成实际计量溯源、版权或专家科学审校。\par
\noindent\textcolor{red!65!black}{[跨辑去重未完成 · 草稿]} \par
\medskip\hrule\medskip
\Needspace{7\baselineskip}
\subsubsection*{UC-ADV9-A86 · 荧光相关光谱FCS停留时间与Stokes–Einstein尺寸反演}
\noindent\textbf{题目}\par
假想球形胶体探针在水中进行3D荧光相关谱测量，横向激光束腰w=0.3000μm，扩散相关时间τD=5.000ms。题设横向FCS关系τD=w²/(4D)；水黏度η=1.000×10\textasciicircum{}-3Pa·s，T=298.15K，Stokes–Einstein D=kBT/(6πηrH)。求扩散系数、流体动力学半径与直径，并估计如果束腰被高估10\%，反演半径会如何偏差。\par
\medskip\noindent\textbf{分步解答}
\begin{enumerate}
\item 束腰转换w=3.000×10\textasciicircum{}-7m，停留时间τ=5×10\textasciicircum{}-3s。
\item FCS横向扩散D=w²/(4τ)=4.5e-12m²/s。
\item Stokes–Einstein反演半径rH=kBT/(6πηD)=48.5293483nm。
\item 对应直径2rH=97.0586967nm。
\item 如果误把束腰增为1.1w，D增为1.21D，推算半径降为原值1/1.21=0.826446281。
\item 真实分子非球形与水合壳导致水动力半径不同于硬球几何半径。
\item 独立用校准式、通量方程或概率守恒重算全部关键结论，归一化残差：0.00e+0、0.00e+0、0.00e+0。
\end{enumerate}
\noindent\textbf{答案摘要}\par
答案1 = 4.5e-12；答案2 = 48.5293483；答案3 = 97.0586967；答案4 = 0.826446281\par
\noindent\textbf{模型前提：}
理想分析模型与题给标准态成立；不代表完成实际计量溯源、版权或专家科学审校。\par
\noindent\textcolor{red!65!black}{[跨辑去重未完成 · 草稿]} \par
\medskip\hrule\medskip
\Needspace{7\baselineskip}
\subsubsection*{UC-ADV9-A87 · ATR-IR渐逝场穿透深度与入射角定量灵敏度}
\noindent\textbf{题目}\par
ATR红外测量用晶体折射率n1=2.400、样品n2=1.400，真空波长λ=6.000μm。题设渐逝场有效穿透深度dp=λ/[2πn1√(sin²θ−(n2/n1)²)]，入射角θ相对于表面法线。分别在45.00°与60.00°测量同一样品，假定弱吸收信号与dp成正比。求两角穿透深度、60°信号相对45°比及发生全反射所需最低入射角，解释靠近临界角时灵敏度增强。\par
\medskip\noindent\textbf{分步解答}
\begin{enumerate}
\item 临界角满足sinθc=n2/n1，故θc=asin(1.4/2.4)=35.6853347°。
\item 45°与60°均超过临界角，公式有实数根。
\item 45°穿透深度dp45=0.995582991μm。
\item 60°穿透深度dp60=0.621606262μm。
\item 假设信号线性于穿透深度，S60/S45=dp60/dp45=0.624364083。
\item 越靠近临界角，根号分母趋小，穿透加深；实际强吸收、色散与有限样品厚度会破坏简单正比。
\item 独立用校准式、通量方程或概率守恒重算全部关键结论，归一化残差：2.22e-16、0.00e+0、0.00e+0。
\end{enumerate}
\noindent\textbf{答案摘要}\par
答案1 = 0.995582991；答案2 = 0.621606262；答案3 = 0.624364083；答案4 = 35.6853347\par
\noindent\textbf{模型前提：}
理想分析模型与题给标准态成立；不代表完成实际计量溯源、版权或专家科学审校。\par
\noindent\textcolor{red!65!black}{[跨辑去重未完成 · 草稿]} \par
\medskip\hrule\medskip
\Needspace{7\baselineskip}
\subsubsection*{UC-ADV9-A88 · 尺寸排阻色谱SEC的对数摩尔质量校准与插值边界}
\noindent\textbf{题目}\par
尺寸排阻色谱以标准聚合物标定log10(M/Da)=a−bVe（Ve单位mL）。两标准分别在Ve1=12.00mL时M1=1.000×10\textasciicircum{}5Da、Ve2=18.00mL时M2=1.000×10\textasciicircum{}4Da。未知聚合物在Ve=15.00mL出峰。求校准截距a、斜率b、未知聚合物推算摩尔质量、若以质量对Ve直接线性插值的错误结果，并指出不同聚合物水动力体积不同导致校准不可普适。\par
\medskip\noindent\textbf{分步解答}
\begin{enumerate}
\item 将两标准取对数，(V1,logM1)=(12,5)、(V2,logM2)=(18,4)。
\item 由两点式b=(5−4)/(18−12)=0.166666667mL\textasciicircum{}-1。
\item 截距a=logM1+bV1=7。
\item 未知logM=a−b×15=4.5，摩尔质量M=31622.7766Da。
\item 若错误对M而不是logM做线性插值，15mL所得质量为55000Da。
\item SEC真正按溶液水动力体积排阻；即便相同分子质量，不同支化和溶剂化也会改变保留。
\item 独立用校准式、通量方程或概率守恒重算全部关键结论，归一化残差：0.00e+0、0.00e+0、0.00e+0。
\end{enumerate}
\noindent\textbf{答案摘要}\par
答案1 = 7；答案2 = 0.166666667；答案3 = 31622.7766；答案4 = 55000\par
\noindent\textbf{模型前提：}
理想分析模型与题给标准态成立；不代表完成实际计量溯源、版权或专家科学审校。\par
\noindent\textcolor{red!65!black}{[跨辑去重未完成 · 草稿]} \par
\medskip\hrule\medskip
\Needspace{7\baselineskip}
\subsubsection*{UC-ADV9-A89 · 差示扫描量热DSC热流峰面积与升温速率的焓换算}
\noindent\textbf{题目}\par
差示扫描量热仪以恒定升温速率β=10.00K/min加热一份质量10.00g的样品。扣除基线后，一次相变的热功率P(T)（W）对温度T积分面积∫P(T)dT=100.0W·K，且无其他相变。求升温速率的SI值、相变吸热J、单位质量焓J/kg，以及错误直接把W·K面积当作J所缺少的因子。\par
\medskip\noindent\textbf{分步解答}
\begin{enumerate}
\item 热流功率P=dQ/dt，恒温速β=dT/dt=10/60=0.166666667K/s。
\item dQ=Pdt=(P/β)dT。
\item 相变热Q=(∫P dT)/β=600J。
\item 样品质量10g=0.010kg，比焓Δh=Q/m=60000J/kg。
\item 若错把100W·K等同100J，漏掉1/β=6s/K的换算因子。
\item 需要检查峰是否已经进行基线扣除与升温速率实际恒定；非恒速应直接按时间积分。
\item 独立用校准式、通量方程或概率守恒重算全部关键结论，归一化残差：0.00e+0、0.00e+0、0.00e+0。
\end{enumerate}
\noindent\textbf{答案摘要}\par
答案1 = 0.166666667；答案2 = 600；答案3 = 60000；答案4 = 6\par
\noindent\textbf{模型前提：}
理想分析模型与题给标准态成立；不代表完成实际计量溯源、版权或专家科学审校。\par
\noindent\textcolor{red!65!black}{[跨辑去重未完成 · 草稿]} \par
\medskip\hrule\medskip
\Needspace{7\baselineskip}
\subsubsection*{UC-ADV9-A90 · 光谱杂散光造成的高吸光非线性压缩与逆校正}
\noindent\textbf{题目}\par
某紫外可见分光仪真实样品吸光度Atrue=2.000，仪器有相对入射光强度1.000\%的杂散光s=0.01000。忽略其它漂移，观测透过率Tobs=(10\textasciicircum{}-Atrue+s)/(1+s)、观测吸光Aobs=−log10(Tobs)。求观测吸光度、低估量、由已知s逆反演真实吸光度，并说明高A区域的杂散光地板和为何不能简单把s直接从Aobs相减。\par
\medskip\noindent\textbf{分步解答}
\begin{enumerate}
\item 真实透过率Ttrue=10\textasciicircum{}-2=0.01000。
\item 实际透过率Tobs=(0.01000+0.01000)/(1.01000)=0.0198019802。
\item 仪器输出Aobs=−log10(Tobs)=1.70329138，低于真实2.000。
\item 吸光低估量Atrue−Aobs=0.296708622。
\item 已知杂散光s可反解Ttrue=(1+s)Tobs−s，恢复Atrue=2。
\item 若Tobs趋于s/(1+s)，仪器出现吸光上限地板；这是透过率通道的加性偏移，不是吸光度加性偏移。
\item 独立用校准式、通量方程或概率守恒重算全部关键结论，归一化残差：0.00e+0、0.00e+0、0.00e+0。
\end{enumerate}
\noindent\textbf{答案摘要}\par
答案1 = 1.70329138；答案2 = 0.296708622；答案3 = 2；答案4 = 0.0198019802\par
\noindent\textbf{模型前提：}
理想分析模型与题给标准态成立；不代表完成实际计量溯源、版权或专家科学审校。\par
\noindent\textcolor{red!65!black}{[跨辑去重未完成 · 草稿]} \par
\medskip\hrule\medskip
\Needspace{7\baselineskip}
\subsubsection*{UC-ADV9-A91 · 放射性示踪活度的化学回收率与衰变双修正}
\noindent\textbf{题目}\par
假想放射性分析物的初始活度A0未知，半衰期t1/2=6.000h。取样后12.00h检测，化学处理对分析物共同回收率r=0.8000，测得回收后活度200.0Bq（两步骤相互独立）。求仅物理衰变剩余比例、总有效测量比例、样品初始活度、如果只补衰变不补回收的错误值，以及初始每秒衰变数。\par
\medskip\noindent\textbf{分步解答}
\begin{enumerate}
\item 放射衰变因子d=2\textasciicircum{}(−t/t1/2)=2\textasciicircum{}−2=0.25。
\item 化学回收和核衰变共同作用，观测比例rd=0.2。
\item 测得Aobs=200Bq，所以初始A0=Aobs/(rd)=1000Bq。
\item 若只做核衰变补偿则得到Awrong=Aobs/d=800Bq，仍低估初始活度。
\item 1Bq定义为每秒一次核衰变，因此初始衰变数率为1000s\textasciicircum{}-1。
\item 如果化学回收率随同位素分馏或衰变子体迁移变化，乘法分离不成立。
\item 独立用校准式、通量方程或概率守恒重算全部关键结论，归一化残差：0.00e+0、0.00e+0、0.00e+0。
\end{enumerate}
\noindent\textbf{答案摘要}\par
答案1 = 0.25；答案2 = 0.2；答案3 = 1000；答案4 = 800\par
\noindent\textbf{模型前提：}
理想分析模型与题给标准态成立；不代表完成实际计量溯源、版权或专家科学审校。\par
\noindent\textcolor{red!65!black}{[跨辑去重未完成 · 草稿]} \par
\medskip\hrule\medskip
\Needspace{7\baselineskip}
\subsubsection*{UC-ADV9-A92 · Job连续变分法对二比一复合物的极大值与化学计量界限}
\noindent\textbf{题目}\par
假想A、B在固定体积、固定总分析浓度Ct=0.006000M下只形成单一A2B配合物，复合物生成非常强，且A、B的自由吸收已扣除，检测信号与[A2B]成正比。改变A的分析摩尔分数xA，证明Job曲线在xA=2/3达到最大，并给出此处A、B各浓度、复合物上界浓度及如果配位平衡常数有限峰形会如何变化。\par
\medskip\noindent\textbf{分步解答}
\begin{enumerate}
\item 给定总浓度Ct，A分析浓度=xCt，B=(1−x)Ct。
\item 形成每mol A2B消耗2mol A与1mol B，故复合物上界c=min(xCt/2,(1−x)Ct)。
\item 最大值发生两项相等xCt/2=(1−x)Ct，求x*=2/3。
\item 此时A浓度=0.004M、B=0.002M。
\item 理想完全结合的复合物上界cmax=Ct/3=0.002M。
\item 若形成常数有限，必须把自由A/B和A2B质量作用联立；曲线变圆滑而不再是分段直线。
\item 独立用校准式、通量方程或概率守恒重算全部关键结论，归一化残差：-4.34e-19、0.00e+0、0.00e+0。
\end{enumerate}
\noindent\textbf{答案摘要}\par
答案1 = 0.666666667；答案2 = 0.004；答案3 = 0.002；答案4 = 0.002\par
\noindent\textbf{模型前提：}
理想分析模型与题给标准态成立；不代表完成实际计量溯源、版权或专家科学审校。\par
\noindent\textcolor{red!65!black}{[跨辑去重未完成 · 草稿]} \par
\medskip\hrule\medskip
\Needspace{7\baselineskip}
\subsubsection*{UC-ADV9-A93 · 含双氯同位素分子M、M+2与M+4的二项分布谱型}
\noindent\textbf{题目}\par
一个假想有机分析物分子每分子恰含两个相互独立的氯原子，天然37Cl原子分数p=0.2420、35Cl分数q=0.7580。忽略13C与其它同位素，37Cl相比35Cl使质谱质量约增加2Da。求分子离子三峰(M、M+2、M+4)的相对概率、相对M峰的比值，以及平均因氯同位素产生的额外质量Da，并讨论不能把M+2峰误解为只含一枚氯。\par
\medskip\noindent\textbf{分步解答}
\begin{enumerate}
\item 两个氯的重同位素个数K服从Binomial(n=2,p=.242)。
\item M峰（零枚37Cl）概率P0=q²=0.574564。
\item M+2峰（恰一枚37Cl）概率P1=2pq=0.366872。
\item M+4峰（两枚37Cl）概率P2=p²=0.058564。
\item 峰概率总和1，M+2相对M的强度比=0.6385224274。
\item 重氯平均个数E[K]=2p，单枚加2Da，故平均质量偏移4p=0.968Da。
\item 独立复核概率总和、数据响应或测量电路约束，数值残差：-1.11e-16、0.00e+0、0.00e+0。
\end{enumerate}
\noindent\textbf{答案摘要}\par
答案1 = 0.574564；答案2 = 0.366872；答案3 = 0.058564；答案4 = 0.638522427；答案5 = 0.968\par
\noindent\textbf{模型前提：}
题设校准、统计与仪器模型按给定条件成立；不代表获得实验计量认证或专家审核。\par
\noindent\textcolor{red!65!black}{[跨辑去重未完成 · 草稿]} \par
\medskip\hrule\medskip
\Needspace{7\baselineskip}
\subsubsection*{UC-ADV9-A94 · 空白噪声与仪器斜率的检出限、定量限及稀释决策}
\noindent\textbf{题目}\par
光谱定量检测的已校准线性斜率s=0.8000吸光度/mM，空白信号标准差σblank=0.006000吸光度，截距已扣。以LOD=3.3σ/s、LOQ=10σ/s定义方法限度。20倍稀释样的净吸光信号S=0.05000。求LOD、LOQ、稀释样及原样的信号逆算浓度，并判断该样本可检出、是否满足定量限；不能把低于LOQ的逆算值自动宣称可靠定量。\par
\medskip\noindent\textbf{分步解答}
\begin{enumerate}
\item LOD=3.3σ/s=0.02475mM。
\item LOQ=10σ/s=0.075mM。
\item 稀释样信号逆算c=S/s=0.0625mM。
\item 原样做20倍稀释，因此形式反推c0=20c=1.25mM。
\item 因LOD<c<LOQ，按题设指标属于可检出而未满足定量限；原样形式浓度不能因此被宣称具有验证过的定量可靠性。
\item 要得到达标定量浓度，需要减少稀释、改善噪声或采用不同校准范围，并检查不确定度。
\item 独立复核概率总和、数据响应或测量电路约束，数值残差：0.00e+0、0.00e+0、0.00e+0。
\end{enumerate}
\noindent\textbf{答案摘要}\par
答案1 = 0.02475；答案2 = 0.075；答案3 = 0.0625；答案4 = 1.25\par
\noindent\textbf{模型前提：}
题设校准、统计与仪器模型按给定条件成立；不代表获得实验计量认证或专家审核。\par
\noindent\textcolor{red!65!black}{[跨辑去重未完成 · 草稿]} \par
\medskip\hrule\medskip
\Needspace{7\baselineskip}
\subsubsection*{UC-ADV9-A95 · Randles阻抗半圆顶点与电荷转移电阻的独立读取}
\noindent\textbf{题目}\par
假想无扩散尾的单一Randles电路，由串联溶液电阻Rs=10.00Ω和并联Rct=200.0Ω、双电层电容Cdl=20.00μF构成。理想Nyquist图为单半圆。求高频/低频实轴截距、半圆顶点频率fmax=1/(2πRctCdl)、顶点复阻抗实部和虚部（约定−ImZ向上）、半圆直径，并说明真实Warburg尾将怎样破坏单半圆假设。\par
\medskip\noindent\textbf{分步解答}
\begin{enumerate}
\item ω→∞时电容趋短路，Zreal→Rs=10Ω。
\item ω→0时电容开路，Zreal→Rs+Rct=210Ω。
\item 顶点满足ω Rct Cdl=1，因此fmax=39.78873577Hz。
\item 顶点实部Rs+Rct/2=110Ω。
\item 顶点虚部ImZ=−Rct/2=-100Ω，Nyquist纵轴若绘制−ImZ则顶点高度为正100Ω。
\item 半圆直径为Rct=200Ω，若出现扩散Warburg或多个时间常数则不再满足单一理想半圆。
\item 独立复核概率总和、数据响应或测量电路约束，数值残差：-1.11e-16、0.00e+0、0.00e+0。
\end{enumerate}
\noindent\textbf{答案摘要}\par
答案1 = 10；答案2 = 210；答案3 = 39.7887358；答案4 = 110；答案5 = -100\par
\noindent\textbf{模型前提：}
题设校准、统计与仪器模型按给定条件成立；不代表获得实验计量认证或专家审核。\par
\noindent\textcolor{red!65!black}{[跨辑去重未完成 · 草稿]} \par
\medskip\hrule\medskip
\Needspace{7\baselineskip}
\subsubsection*{UC-ADV9-A96 · 手性旋光度、比旋光与两对映体质量分数反演}
\noindent\textbf{题目}\par
在选定波长温度下，纯R对映体比旋光[α]=+80.00°·dm\textasciicircum{}-1·(g/mL)\textasciicircum{}-1，纯S为−80.00。外消旋以外未知混合物总浓度0.2000g/mL、比色管光程1.000dm，实测旋光角+8.000°，忽略其它旋光杂质。求R优势对映体过量ee、R/S各自质量浓度、R/S质量分数以及若错误把观察旋光角直接除纯品比旋光而忽略光程/浓度产生的误判。\par
\medskip\noindent\textbf{分步解答}
\begin{enumerate}
\item 混合旋光角αobs=[α]l(cR−cS)。
\item 总质量浓度cR+cS=0.2000g/mL。
\item ee=(cR−cS)/(cR+cS)=αobs/([α]lC)=0.5。
\item R浓度cR=C(1+ee)/2=0.15g/mL。
\item S浓度cS=C(1−ee)/2=0.05g/mL，R质量分数=0.75。
\item 仅旋光测量无法证明无其它旋光杂质，条件和已知纯品比旋光是反演前提。
\item 独立复核概率总和、数据响应或测量电路约束，数值残差：1.78e-15、0.00e+0、-1.11e-16。
\end{enumerate}
\noindent\textbf{答案摘要}\par
答案1 = 0.5；答案2 = 0.15；答案3 = 0.05；答案4 = 0.75\par
\noindent\textbf{模型前提：}
题设校准、统计与仪器模型按给定条件成立；不代表获得实验计量认证或专家审核。\par
\noindent\textcolor{red!65!black}{[跨辑去重未完成 · 草稿]} \par
\medskip\hrule\medskip
\Needspace{7\baselineskip}
\subsubsection*{UC-ADV9-A97 · 荧光偏振Perrin方程与旋转相关时间推断半径}
\noindent\textbf{题目}\par
球形荧光探针的基本各向异性r0=0.4000，实测稳态各向异性r=0.2500，激发态寿命τ=4.000ns。题设Perrin式r=r0/(1+τ/θ)，并假设Stokes黏性旋转θ=ηV/(kBT)，水黏度η=1.000×10\textasciicircum{}-3Pa·s，T=298.15K。求旋转相关时间θ、旋转扩散常数Drot=1/(6θ)、等效球体体积与半径，说明非球形形状和标记局部转动会造成偏差。\par
\medskip\noindent\textbf{分步解答}
\begin{enumerate}
\item 由r0/r=1+τ/θ解得θ=τ/(r0/r−1)=6.666666667ns。
\item 旋转扩散常数Drot=1/(6θ)=25000000s\textasciicircum{}-1。
\item 在理想球形黏性模型中等效体积V=kBT θ/η=2.744269996e-26m³。
\item 球形体积V=4πa³/3，半径a=[3V/(4π)]\textasciicircum{}(1/3)=1.871167832nm。
\item 将θ逆代Perrin式恢复r=0.2500；若探针自由转动独立于载体，推得的半径可能不是分子整体半径。
\item 需要同时测得寿命和各向异性，单靠稳态各向异性不能分离寿命和转动时间。
\item 独立复核概率总和、数据响应或测量电路约束，数值残差：0.00e+0、0.00e+0、0.00e+0。
\end{enumerate}
\noindent\textbf{答案摘要}\par
答案1 = 6.66666667；答案2 = 25000000；答案3 = 2.74427e-26；答案4 = 1.87116783\par
\noindent\textbf{模型前提：}
题设校准、统计与仪器模型按给定条件成立；不代表获得实验计量认证或专家审核。\par
\noindent\textcolor{red!65!black}{[跨辑去重未完成 · 草稿]} \par
\medskip\hrule\medskip
\Needspace{7\baselineskip}
\subsubsection*{UC-ADV9-A98 · 高分辨质谱质量准确度ppm与峰分辨率的区别}
\noindent\textbf{题目}\par
高分辨质谱某已知离子理论m/z=500.0000Da，实测峰中心500.0125Da，峰的FWHM=0.02500Da。另一个相邻峰与本峰中心相差0.03300Da，假设二者峰宽相等。求绝对质量偏差、ppm偏差、质量分辨率R=m/FWHM以及峰间距/FWHM之比，并说明准确度与分辨率是不同性能指标、仅FWHM不能保证基线分离。\par
\medskip\noindent\textbf{分步解答}
\begin{enumerate}
\item 质量准确度偏差Δm=500.0125−500.0000=0.0125Da。
\item ppm=(Δm/mtheory)×10\textasciicircum{}6=25ppm。
\item 分辨率R=m/FWHM=20000。
\item 两相邻峰中心间距与FWHM之比=1.32。
\item R大不代表峰中心完全准确；质量偏差取决于校准而非半峰宽。
\item 峰基线是否分离还与峰形、强度比及噪声有关，不能单靠中心间距大于FWHM保证。
\item 独立复核概率总和、数据响应或测量电路约束，数值残差：0.00e+0、0.00e+0、0.00e+0。
\end{enumerate}
\noindent\textbf{答案摘要}\par
答案1 = 0.0125；答案2 = 25；答案3 = 20000；答案4 = 1.32\par
\noindent\textbf{模型前提：}
题设校准、统计与仪器模型按给定条件成立；不代表获得实验计量认证或专家审核。\par
\noindent\textcolor{red!65!black}{[跨辑去重未完成 · 草稿]} \par
\medskip\hrule\medskip
\Needspace{7\baselineskip}
\subsubsection*{UC-ADV9-A99 · Nikolskii–Eisenman离子选择电极的二价干扰校正}
\noindent\textbf{题目}\par
理想A+离子选择性电极同时受到B2+离子干扰，标准态活度均以1mol/L归一为无量纲，干扰系数KA,B=0.02000。Nikolskii–Eisenman式E=E0+(RT/F)ln(aA+KA,B aB\textasciicircum{}(1/2))，298.15K E0=0.2000V。给定aA=0.001000、aB=0.0004000。求有效等效活度、测得电势、干扰项占等效活度的分数，以及忽略B所预测的电势误差，说明为何二价干扰项要取平方根。\par
\medskip\noindent\textbf{分步解答}
\begin{enumerate}
\item 在参考标准态下aA与aB均无量纲，干扰项K aB\textasciicircum{}(1/2)=0.0004。
\item 总等效活度aeff=aA+K√aB=0.0014。
\item 由题给式E=E0+(RT/F)ln(aeff)=0.0311667905V。
\item B干扰占等效信号活度分数=0.2857142857。
\item 忽略B时错误电势Ewrong=E0+(RT/F)ln(aA)=0.02252195094V，差值=0.008644839562V。
\item 二价干扰离子指数zA/zB=1/2来自不同价态对电化学势的影响，是模型约定，真实选择性系数需实验标定。
\item 独立复核概率总和、数据响应或测量电路约束，数值残差：-5.42e-20、0.00e+0、5.42e-20。
\end{enumerate}
\noindent\textbf{答案摘要}\par
答案1 = 0.0014；答案2 = 0.0311667905；答案3 = 0.285714286；答案4 = 0.00864483956\par
\noindent\textbf{模型前提：}
题设校准、统计与仪器模型按给定条件成立；不代表获得实验计量认证或专家审核。\par
\noindent\textcolor{red!65!black}{[跨辑去重未完成 · 草稿]} \par
\medskip\hrule\medskip
\Needspace{7\baselineskip}
\subsubsection*{UC-ADV9-A100 · 带高斯先验的线性测量浓度Bayes后验与不确定度}
\noindent\textbf{题目}\par
假想待测分析物浓度c单位mM，测量前已有独立先验c\textasciitilde{}N(1.000,0.4000²)，仪器线性模型y=2.000c+0.2000+ε，ε\textasciitilde{}N(0,0.2000²)，测得y=2.600。假定校准斜率和截距确定、先验与误差独立。求单凭测量的c估计及σc、合并后验均值、标准差和近似95\%可信区间，并说明这个Bayes区间依赖先验与正态误差假设。\par
\medskip\noindent\textbf{分步解答}
\begin{enumerate}
\item 信号反演的观测浓度ycor=(y−0.2)/2=1.2mM。
\item 观测噪声换算为浓度σy=0.2000/2=0.1mM。
\item 先验精度P0=1/0.4²=6.25，观测精度Py=1/0.1²=100。
\item 高斯后验均值μ=(P0μ0+Pyμy)/(P0+Py)=1.188235294mM。
\item 后验标准差σpost=[P0+Py]\textasciicircum{}(-1/2)=0.09701425001mM。
\item 近似95\%后验区间μ±1.96σ=[0.9980873641,1.378383224]mM；不是脱离先验的频率学置信区间。
\item 独立复核概率总和、数据响应或测量电路约束，数值残差：1.42e-14、0.00e+0、0.00e+0。
\end{enumerate}
\noindent\textbf{答案摘要}\par
答案1 = 1.2；答案2 = 0.1；答案3 = 1.18823529；答案4 = 0.09701425；答案5 = 0.998087364；答案6 = 1.37838322\par
\noindent\textbf{模型前提：}
题设校准、统计与仪器模型按给定条件成立；不代表获得实验计量认证或专家审核。\par
\noindent\textcolor{red!65!black}{[跨辑去重未完成 · 草稿]} \par
\medskip\hrule\medskip
\subsection*{有机化学}
\addcontentsline{toc}{subsection}{第9辑 / 有机化学}
\Needspace{7\baselineskip}
\subsubsection*{UC-ADV9-O73 · E2消除反应反式共平面构象占据与产物速率选择性}
\noindent\textbf{题目}\par
假想可自由互变的卤代环己烷有anti-periplanar消除可达构象A与非anti构象S，A比S的自由能高2.000kJ/mol，温度298.15K。快速构象平衡，A消除一级kA=0.8000min\textasciicircum{}-1，S的另一竞争反应一级kS=0.04000min\textasciicircum{}-1，二路径给不同产物。求A的构象分数、整体消耗kobs、anti消除产物的初速分数及不正确地以最稳定构象直接预测主要产物会产生的概念错误。\par
\medskip\noindent\textbf{分步解答}
\begin{enumerate}
\item A对S的Boltzmann比pA/pS=exp(−ΔG/RT)。
\item 由pA+pS=1解pA=0.3085747544、pS=0.6914252456。
\item 消耗有效一级常数kobs=pA kA+pS kS=0.2745168133min\textasciicircum{}-1。
\item anti消除分支贡献的产物分数fA=(pAkA)/kobs=0.8992520367。
\item 尽管A占据较少，因其内禀消除快，产物仍可由A主导，说明需按构象概率乘路径速率。
\item 若构象互变不能快于消耗，固定Boltzmann加权不成立，必须解双态耦合动力学。
\item 独立重算势垒、速率与成键反应守恒或对称性计数，数值残差：0.00e+0、1.39e-17、0.00e+0。
\end{enumerate}
\noindent\textbf{答案摘要}\par
答案1 = 0.308574754；答案2 = 0.274516813；答案3 = 0.899252037；答案4 = 0.691425246\par
\noindent\textbf{模型前提：}
假想的选择性或动力学模型满足题设，真实机理与模型新颖性仍待专家审查。\par
\noindent\textcolor{red!65!black}{[跨辑去重未完成 · 草稿]} \par
\medskip\hrule\medskip
\Needspace{7\baselineskip}
\subsubsection*{UC-ADV9-O74 · SN1与SN2平行反应的反应级数和速率常数辨识}
\noindent\textbf{题目}\par
假想溶剂解离SN1与双分子SN2同时消耗同一底物S，分别以一级常数k1和二级常数k2作用于固定过量亲核试剂Nu，故kobs=k1+k2[Nu]。独立测得[Nu]=0.1000M时kobs=0.04000min\textasciicircum{}-1，[Nu]=0.3000M时kobs=0.08000min\textasciicircum{}-1。求k1、k2、[Nu]=0.2000M时的kobs及SN2产品份额，说明只测单个Nu浓度不能分离两个路径。\par
\medskip\noindent\textbf{分步解答}
\begin{enumerate}
\item 用两条件作差kobs2−kobs1=k2(Nu2−Nu1)。
\item 由数据得二级常数k2=0.2L mol\textasciicircum{}-1min\textasciicircum{}-1。
\item 代入任一点得到一级常数k1=0.02min\textasciicircum{}-1。
\item 当Nu=0.2000M时kobs=k1+k2Nu=0.06min\textasciicircum{}-1。
\item SN2路径速率分数fSN2=k2Nu/kobs=0.6666666667。
\item 只凭一个Nu条件测到的总一级常数对k1/k2有无限多组解，需要多个亲核试剂浓度测量。
\item 独立重算势垒、速率与成键反应守恒或对称性计数，数值残差：0.00e+0、-1.39e-17、0.00e+0。
\end{enumerate}
\noindent\textbf{答案摘要}\par
答案1 = 0.02；答案2 = 0.2；答案3 = 0.06；答案4 = 0.666666667\par
\noindent\textbf{模型前提：}
假想的选择性或动力学模型满足题设，真实机理与模型新颖性仍待专家审查。\par
\noindent\textcolor{red!65!black}{[跨辑去重未完成 · 草稿]} \par
\medskip\hrule\medskip
\Needspace{7\baselineskip}
\subsubsection*{UC-ADV9-O75 · Hammett取代基常数与并行反应过渡态敏感性}
\noindent\textbf{题目}\par
某假想芳香取代反应在298.15K满足Hammett关系log10(kX/kH)=ρσ。测得带标准Hammett常数σX=+0.3000的取代基使kX/kH=2.000；另一取代基σY=−0.2000。求反应常数ρ、Y的预测相对速率、Y相对H过渡态活化自由能差ΔΔG‡=−RTln(kY/kH)和两个取代基的速率倍率之比，解释此线性经验式只在同反应机理系列内适用。\par
\medskip\noindent\textbf{分步解答}
\begin{enumerate}
\item 由X基团数据ρ=log10(2)/0.3000=1.003433319。
\item Y的Hammett预测log10(kY/kH)=ρσY。
\item 所以kY/kH=10\textasciicircum{}(ρσY)=0.6299605249。
\item 活化自由能差ΔΔG‡=−RTln(kY/kH)=1.145521384kJ/mol。
\item X对Y的速率比(kX/kH)/(kY/kH)=3.174802104。
\item Hammett斜率不证明所有取代基都同机理；若发生酸碱、共振或反应决速步变化，必须重新检验线性关系。
\item 独立重算势垒、速率与成键反应守恒或对称性计数，数值残差：0.00e+0、2.78e-17、0.00e+0。
\end{enumerate}
\noindent\textbf{答案摘要}\par
答案1 = 1.00343332；答案2 = 0.629960525；答案3 = 1.14552138；答案4 = 3.1748021\par
\noindent\textbf{模型前提：}
假想的选择性或动力学模型满足题设，真实机理与模型新颖性仍待专家审查。\par
\noindent\textcolor{red!65!black}{[跨辑去重未完成 · 草稿]} \par
\medskip\hrule\medskip
\Needspace{7\baselineskip}
\subsubsection*{UC-ADV9-O76 · 分子内内酯化有效摩尔浓度与分子间竞争反应}
\noindent\textbf{题目}\par
一份假想羟基酸底物可通过一级分子内内酯化生成内酯P（kintra=0.04000s\textasciicircum{}-1），也可与外部亲核剂B发生二级分子间竞争反应生成Q（kinter=1.000L mol\textasciicircum{}-1s\textasciicircum{}-1）。B由大库维持自由浓度0.03000M，产物不反应。求两个有效一级常数、P产物比例、底物消耗半衰期及有效摩尔浓度EM=kintra/kinter，并解释EM是速率比较量而不是体系真实化学浓度。\par
\medskip\noindent\textbf{分步解答}
\begin{enumerate}
\item 分子间通道伪一级常数kinter,obs=kinter[B]=0.03s\textasciicircum{}-1。
\item 总消耗常数ktot=kintra+kinter,obs=0.07s\textasciicircum{}-1。
\item P累计产物分数kintra/ktot=0.5714285714。
\item 底物半衰期t1/2=ln2/ktot=9.902102579s。
\item 有效摩尔浓度EM=kintra/kinter=0.04M，它使两路径的速率竞争在[B]=EM时相等。
\item 当B不由外库维持时，第二通道会消耗B，总一级近似将失效，必须额外写B物料守恒。
\item 独立重算势垒、速率与成键反应守恒或对称性计数，数值残差：6.94e-18、0.00e+0、0.00e+0。
\end{enumerate}
\noindent\textbf{答案摘要}\par
答案1 = 0.03；答案2 = 0.07；答案3 = 0.571428571；答案4 = 9.90210258；答案5 = 0.04\par
\noindent\textbf{模型前提：}
假想的选择性或动力学模型满足题设，真实机理与模型新颖性仍待专家审查。\par
\noindent\textcolor{red!65!black}{[跨辑去重未完成 · 草稿]} \par
\medskip\hrule\medskip
\Needspace{7\baselineskip}
\subsubsection*{UC-ADV9-O77 · 光引发自由基链反应的量子产率与平均链增长长度}
\noindent\textbf{题目}\par
假想光致自由基链反应在450.0nm波长实验中有效吸收光子量为2.000×10\textasciicircum{}-5mol photons，实验产物形成1.000×10\textasciicircum{}-3mol，光化学引发效率为0.8000条新链/mol被吸收光子，所有引发链都能产生题列产品。求表观产品量子产率、每根有效引发链的平均产品数、吸收光子总能量以及如果简单认为一个光子至多形成一个产品将出现的物料矛盾。\par
\medskip\noindent\textbf{分步解答}
\begin{enumerate}
\item 产品量子产率Φ=nP/nphoton=50mol/mol。
\item 真正有效引发链条数nchain=η nphoton=0.000016mol。
\item 每根链对应平均产品数L=nP/nchain=62.5。
\item 光子总能量Eabs=nphoton NA hc/λ=5.316736251J。
\item Φ≫1可以由链式增长解释，并不违背光能量守恒，因为化学反应可另释放或消耗自由能。
\item 如果链并不都存活或每链有终止副反应，则推断的平均链长需按实际引发与终止效率再校正。
\item 独立重算势垒、速率与成键反应守恒或对称性计数，数值残差：0.00e+0、0.00e+0、0.00e+0。
\end{enumerate}
\noindent\textbf{答案摘要}\par
答案1 = 50；答案2 = 62.5；答案3 = 5.31673625；答案4 = 1.6e-05\par
\noindent\textbf{模型前提：}
假想的选择性或动力学模型满足题设，真实机理与模型新颖性仍待专家审查。\par
\noindent\textcolor{red!65!black}{[跨辑去重未完成 · 草稿]} \par
\medskip\hrule\medskip
\Needspace{7\baselineskip}
\subsubsection*{UC-ADV9-O78 · 单取代苯亲电芳香取代的等价位点数与区域选择性}
\noindent\textbf{题目}\par
某假想单取代苯亲电芳香取代中，分子有2个对称等价ortho位置、2个meta位置、1个para位置，且在相同初始反应条件下每个位置的内禀反应速率分别ko=0.8000、km=0.1000、kp=1.600（共同任意速率单位）。假设一次取代且不存在其他路径。求三类产物累计初速比例、ortho/para比、产物meta占比和忽略位点简并度的选择性预测误差，说明“每位点反应性”与“总体产物份额”不同。\par
\medskip\noindent\textbf{分步解答}
\begin{enumerate}
\item 先计算各位点类别的总反应权重Wortho=2ko=1.6。
\item Wmeta=2km=0.2，Wpara=kp=1.6。
\item 总权重Wtot=3.4。
\item 三类产物分数分别为0.4705882353、0.05882352941、0.4705882353。
\item ortho/para产物比=(2ko)/kp=1，虽每个ortho位点内禀较慢，两个位置使总量与para可能相当。
\item 若错误按ko:km:kp三数直接归一，则meta分数被错估为0.04而非0.05882352941。
\item 独立重算势垒、速率与成键反应守恒或对称性计数，数值残差：0.00e+0、0.00e+0、0.00e+0。
\end{enumerate}
\noindent\textbf{答案摘要}\par
答案1 = 0.470588235；答案2 = 0.0588235294；答案3 = 0.470588235；答案4 = 1；答案5 = 0.0588235294\par
\noindent\textbf{模型前提：}
假想的选择性或动力学模型满足题设，真实机理与模型新颖性仍待专家审查。\par
\noindent\textcolor{red!65!black}{[跨辑去重未完成 · 草稿]} \par
\medskip\hrule\medskip
\Needspace{7\baselineskip}
\subsubsection*{UC-ADV9-O79 · Hückel六元芳香环的π能量稳定化与乙烯局域参照}
\noindent\textbf{题目}\par
简单Hückel近似中苯六个π电子占据一条α+2β能级（2电子）和两条α+β能级（共4电子），β=−75.00kJ/mol；三分子互不共轭乙烯的六个π电子每分子只占据α+β（共6电子）。求苯与三乙烯的π电子总能表达式、苯相对局域化乙烯的额外稳定化能量和每个π电子的平均稳定化贡献，并解释Hückel稳定化只是一种比较参照，不等于真实苯的实验共振能。\par
\medskip\noindent\textbf{分步解答}
\begin{enumerate}
\item 苯π能量Ebenz=2(α+2β)+4(α+β)=6α+8β。
\item 三独立乙烯参考Eref=6(α+β)=6α+6β。
\item 两者能量差ΔE=Ebenz−Eref=2β=-150kJ/mol。
\item β取负值，ΔE为负，因此苯相对局域参照具有更低π能。
\item 以绝对值表示额外稳定化为150kJ/mol苯。
\item 平均每个π电子的额外能量差ΔE/6=-25kJ/mol π电子，真实化学反应能还依赖σ骨架、几何形变及电子相关。
\item 独立重算势垒、速率与成键反应守恒或对称性计数，数值残差：0.00e+0、0.00e+0、0.00e+0。
\end{enumerate}
\noindent\textbf{答案摘要}\par
答案1 = -150；答案2 = 150；答案3 = -25\par
\noindent\textbf{模型前提：}
假想的选择性或动力学模型满足题设，真实机理与模型新颖性仍待专家审查。\par
\noindent\textcolor{red!65!black}{[跨辑去重未完成 · 草稿]} \par
\medskip\hrule\medskip
\Needspace{7\baselineskip}
\subsubsection*{UC-ADV9-O80 · Karplus vicinal氢核耦合常数的二解与NOE构象排除}
\noindent\textbf{题目}\par
某有机分子邻近两质子的vicinal NMR耦合常数为3J=5.000Hz，采用经验Karplus式3J=Acos²φ+Bcosφ+C，A=8.000Hz、B=−1.000Hz、C=1.000Hz，φ∈[0,180°]。求代数解cosφ、两个可能二面角，并结合独立NOE判断已知φ>90°时选哪个构象，说明单个J值不足以唯一确定三维构象。\par
\medskip\noindent\textbf{分步解答}
\begin{enumerate}
\item 令x=cosφ，则8x²−x+1=5，整理8x²−x−4=0。
\item 判别式Δ=1+128=129，两个实根x±=(1±√129)/16。
\item 两个允许的二面角φ1=acos(x+)=39.43339001°、φ2=acos(x−)=130.343118°。
\item 因独立NOE要求φ>90°，应选择φ2=130.343118°。
\item 把两个角的cos分别代回Karplus式均得5.000Hz，表明单一J值有二解歧义。
\item 实际Karplus参数随取代基与溶剂变化，NOE也需结合距离和分子动力学判断。
\item 独立重算势垒、速率与成键反应守恒或对称性计数，数值残差：0.00e+0、0.00e+0、4.44e-16。
\end{enumerate}
\noindent\textbf{答案摘要}\par
答案1 = 0.772363543；答案2 = -0.647363543；答案3 = 39.43339；答案4 = 130.343118\par
\noindent\textbf{模型前提：}
假想的选择性或动力学模型满足题设，真实机理与模型新颖性仍待专家审查。\par
\noindent\textcolor{red!65!black}{[跨辑去重未完成 · 草稿]} \par
\medskip\hrule\medskip
\Needspace{7\baselineskip}
\subsubsection*{UC-ADV9-O81 · 高分辨质谱式量、不饱和度与氮规则互相约束}
\noindent\textbf{题目}\par
一个未知有机分子的候选名义分子式为C10H11NO2，不含卤素，假定中性闭壳层。名义原子质量C=12、H=1、N=14、O=16，给定不饱和度DBE=(2C+2+N−H)/2。若13C自然丰度约1.100\%、忽略其它重同位素，求名义质量、DBE、M+1/M峰一阶近似比例及氮规则是否相容，并说明为什么同一不饱和度可能存在许多结构异构体。\par
\medskip\noindent\textbf{分步解答}
\begin{enumerate}
\item 按式量得M=10×12+11+14+32=177Da。
\item DBE=(2×10+2+1−11)/2=6。
\item 分子有奇数氮，名义分子离子质量应为奇数；177符合氮规则。
\item 单个13C替代一阶峰比例为10×1.1\%=11\%。
\item DBE表示环数与π键当量之和，不能仅凭DBE确定具体骨架。
\item 真实HRMS应使用精确原子质量、加合状态与完整同位素分布，而非名义质量。
\item 复核：按原始守恒、电子计数及热力学关系回代，数值误差：0.00e+0、0.00e+0、0.00e+0。
\end{enumerate}
\noindent\textbf{答案摘要}\par
答案1 = 177；答案2 = 6；答案3 = 0.11；答案4 = 11\par
\noindent\textbf{模型前提：}
题设为理想化有机化学模型；模型新颖性、实际机理与专家审定尚未完成。\par
\noindent\textcolor{red!65!black}{[跨辑去重未完成 · 草稿]} \par
\medskip\hrule\medskip
\Needspace{7\baselineskip}
\subsubsection*{UC-ADV9-O82 · 铁羰基光致脱配位与氢氧化加成的电子计数}
\noindent\textbf{题目}\par
一个假想配位循环从Fe(CO)5出发，光照脱去一个CO生成Fe(CO)4，再与H2氧化加成形成Fe(H)2(CO)4。规定Fe(0)按d8、Fe(II)按d6计，每个CO在中性计数法贡献2电子、氢化物在离子计数法各给2电子。求三阶段金属电子总数、铁形式氧化态和16电子中间体距18电子的差额；不得在加成后继续把Fe(II)误按d8计算。\par
\medskip\noindent\textbf{分步解答}
\begin{enumerate}
\item Fe(CO)5的Fe(0)为d8，5个CO共给10电子，总数18。
\item 失去一个CO仍为Fe(0)，Fe(CO)4的电子数8+8=16。
\item 氧化加成后两个氢化物各−I，铁形式氧化态从0提高到+II。
\item Fe(II)是d6，终态电子计数6+4×2+2×2=18。
\item Fe(CO)4距18e规则有2e缺额；终态重新达到18e。
\item 电子计数并不证明光解和氧化加成在真实体系中必然沿这一唯一机制进行。
\item 复核：按原始守恒、电子计数及热力学关系回代，数值误差：0.00e+0、0.00e+0、0.00e+0。
\end{enumerate}
\noindent\textbf{答案摘要}\par
答案1 = 18；答案2 = 16；答案3 = 18；答案4 = 2\par
\noindent\textbf{模型前提：}
题设为理想化有机化学模型；模型新颖性、实际机理与专家审定尚未完成。\par
\noindent\textcolor{red!65!black}{[跨辑去重未完成 · 草稿]} \par
\medskip\hrule\medskip
\Needspace{7\baselineskip}
\subsubsection*{UC-ADV9-O83 · 环状单体聚合天花板温度及跨温度平衡方向}
\noindent\textbf{题目}\par
假设一份有机环状单体每重复单元开环聚合的标准反应焓ΔH=−35.00kJ/mol、标准熵ΔS=−110.0J/(mol K)，在298.15至340K之间近似恒定。计算标准反应自由能变为零的天花板温度、两个温度下的ΔG与标准态平衡常数K=exp(−ΔG/RT)，判断升温后聚合与解聚的有利方向，并解释实际单体活度会如何影响可达到的平衡转化。\par
\medskip\noindent\textbf{分步解答}
\begin{enumerate}
\item 用ΔG(T)=ΔH−TΔS求相平衡热力学驱动力。
\item 令ΔG(Tc)=0，得到Tc=ΔH/ΔS=318.181818K。
\item 298.15K时ΔG=-2203.5J/mol，K=2.43240842。
\item 340K时ΔG=2400J/mol，K=0.427850549。
\item 聚合熵变为负，升温使聚合自由能增大，在天花板温度以上标准态下解聚有利。
\item 实际反应商需纳入剩余单体与链端活动度，不可将标准K直接等同最终转化率。
\item 复核：按原始守恒、电子计数及热力学关系回代，数值误差：0.00e+0、-1.11e-16、0.00e+0。
\end{enumerate}
\noindent\textbf{答案摘要}\par
答案1 = 318.181818；答案2 = -2203.5；答案3 = 2400；答案4 = 2.43240842；答案5 = 0.427850549\par
\noindent\textbf{模型前提：}
题设为理想化有机化学模型；模型新颖性、实际机理与专家审定尚未完成。\par
\noindent\textcolor{red!65!black}{[跨辑去重未完成 · 草稿]} \par
\medskip\hrule\medskip
\Needspace{7\baselineskip}
\subsubsection*{UC-ADV9-O84 · 固相五步肽链偶联的累积完整率与单缺失链分布}
\noindent\textbf{题目}\par
假想固相肽合成从0.4000mmol树脂有效位点出发，完成目标肽需要五步独立偶联，每步成功率均为0.9700。失败链继续参加后续反应并产生缺失肽，最终裂解和分离对各链具有统一回收率0.9000。求五步全部偶联成功的链比例、恰有一次失败的链比例、实际得到完整肽mmol和误把一次偶联率当作总体偶联率所造成的过高估计倍数。\par
\medskip\noindent\textbf{分步解答}
\begin{enumerate}
\item 五步都成功概率Pfull=0.97\textasciicircum{}5=0.858734026。
\item 恰一次失败可出现在五个步骤中的任意一步，概率P1=5×0.03×0.97\textasciicircum{}4=0.132793922。
\item 裂解前完整链物质的量为树脂载量乘Pfull=0.34349361mmol。
\item 回收后完整肽量为0.4000×Pfull×0.9000=0.309144249mmol。
\item 若仅使用单步成功率0.97，则完整肽量被高估倍数0.97/Pfull=1.12956977。
\item 真实序列可能影响各步反应性与裂解率，本题按独立同分布概率作为题设。
\item 复核：按原始守恒、电子计数及热力学关系回代，数值误差：0.00e+0、0.00e+0、0.00e+0。
\end{enumerate}
\noindent\textbf{答案摘要}\par
答案1 = 0.858734026；答案2 = 0.132793922；答案3 = 0.309144249；答案4 = 1.12956977\par
\noindent\textbf{模型前提：}
题设为理想化有机化学模型；模型新颖性、实际机理与专家审定尚未完成。\par
\noindent\textcolor{red!65!black}{[跨辑去重未完成 · 草稿]} \par
\medskip\hrule\medskip
\Needspace{7\baselineskip}
\subsubsection*{UC-ADV9-O85 · Suzuki交叉偶联与双底物副反应的限制收率审计}
\noindent\textbf{题目}\par
假想Suzuki交叉偶联以ArBr=0.01000mol、芳基硼酸=0.01200mol为初始原料，目标产物每mol各消耗1mol底物。另有脱卤副反应消耗0.001000mol ArBr，硼酸自偶联与氧化合计消耗0.004000mol硼酸，之后偶联只使剩余限制底物的75.00\%转为目标产物。计算剩余限制底物、产物mol、两个剩余底物mol及目标产物对初始ArBr的收率，并列出两个独立元素物料账本。\par
\medskip\noindent\textbf{分步解答}
\begin{enumerate}
\item 扣除副反应后可用于偶联的ArBr为0.009000mol。
\item 硼酸剩余可供偶联为0.008000mol，配比1:1，因此理论最多生成目标产品0.008mol。
\item 偶联实际生成P=0.75×lim=0.006mol。
\item 剩余ArBr=0.010−0.001−P=0.003mol。
\item 剩余硼酸=0.012−0.004−P=0.002mol，两个元素账本分别闭合。
\item 相对原始ArBr的产物收率P/0.010=0.6，不是75.00\%。
\item 独立核验全部反应物账本、状态概率和动力学物理解残差：1.73e-18、0.00e+0、0.00e+0。
\end{enumerate}
\noindent\textbf{答案摘要}\par
答案1 = 0.008；答案2 = 0.006；答案3 = 0.003；答案4 = 0.002；答案5 = 0.6\par
\noindent\textbf{模型前提：}
仅假设题给真实物料与动力学机制；跨辑去重、实验真实性及人类专家核定尚待完成。\par
\noindent\textcolor{red!65!black}{[跨辑去重未完成 · 草稿]} \par
\medskip\hrule\medskip
\Needspace{7\baselineskip}
\subsubsection*{UC-ADV9-O86 · 立体异构产物由初始动力学选择性趋向热平衡的量化}
\noindent\textbf{题目}\par
一份有机反应在停止生产后留下A/B立体异构产物初始分数0.9000/0.1000。随后恒温条件下两异构体以A→B一级kAB=0.1000h\textasciicircum{}-1、B→A一级kBA=0.2000h\textasciicircum{}-1互变，物料总量保持1mol。求10.00h后A/B、最终热力学平衡A比例、初始与10h的A/B比和相对初始高选择性的减少量，并说明最终产物比例可能与原生成过渡态的动力学选择性无关。\par
\medskip\noindent\textbf{分步解答}
\begin{enumerate}
\item 定义A占比a，b=1−a，动力学da/dt=−kAB a+kBA(1−a)。
\item 平衡A分数aeq=kBA/(kAB+kBA)=0.6666666667。
\item 从a0=.9解a(t)=aeq+(a0−aeq)e\textasciicircum{}−(kAB+kBA)t。
\item 10h后A分数=0.6782836493，B分数=0.3217163507。
\item 起始A/B=9，10h比=2.108328183，因可逆互变趋于aeq/(1−aeq)。
\item 高动力学初产比例在后续异构化中会衰减，因此须考虑取样与快速淬灭。
\item 独立核验全部反应物账本、状态概率和动力学物理解残差：0.00e+0、1.39e-17、-3.64e-17。
\end{enumerate}
\noindent\textbf{答案摘要}\par
答案1 = 0.678283649；答案2 = 0.321716351；答案3 = 0.666666667；答案4 = 2.10832818；答案5 = 0.221716351\par
\noindent\textbf{模型前提：}
仅假设题给真实物料与动力学机制；跨辑去重、实验真实性及人类专家核定尚待完成。\par
\noindent\textcolor{red!65!black}{[跨辑去重未完成 · 草稿]} \par
\medskip\hrule\medskip
\Needspace{7\baselineskip}
\subsubsection*{UC-ADV9-O87 · 等规与间规序列Bernoulli立构三联体概率}
\noindent\textbf{题目}\par
某理想有机聚合链每个相邻单体立体关系独立地以概率p=0.7000形成meso二联体m，以概率1−p形成racemo二联体r。三个连续单体包含相邻两个二联体关系。求mm、mr/rm合计、rr三联体分数，连续m关系段首次出现r之前的平均m段长度（包含起始m），以及若误用0.7作为mm概率的偏差；说明真实催化立构聚合可能有高阶记忆。\par
\medskip\noindent\textbf{分步解答}
\begin{enumerate}
\item 两个连续m的概率Pmm=p²=0.49。
\item 恰一m一r共有mr与rm两条序列，合计概率2p(1−p)=0.42。
\item 两个r关系概率Prr=(1−p)²=0.09。
\item 三类概率和1。
\item 连续meso段的几何分布以首次r终止，平均m个数按此长度定义为1/(1−p)=3.333333333。
\item 错误以单二联体p当三联体Pmm会高估p−p²=0.21，高阶Markov链需独立测量。
\item 独立核验全部反应物账本、状态概率和动力学物理解残差：0.00e+0、0.00e+0、0.00e+0。
\end{enumerate}
\noindent\textbf{答案摘要}\par
答案1 = 0.49；答案2 = 0.42；答案3 = 0.09；答案4 = 3.33333333；答案5 = 0.21\par
\noindent\textbf{模型前提：}
仅假设题给真实物料与动力学机制；跨辑去重、实验真实性及人类专家核定尚待完成。\par
\noindent\textcolor{red!65!black}{[跨辑去重未完成 · 草稿]} \par
\medskip\hrule\medskip
\Needspace{7\baselineskip}
\subsubsection*{UC-ADV9-O88 · 力致有机小环开环的Bell势垒降低与半衰期}
\noindent\textbf{题目}\par
假想单分子机械力致开环服从Bell模型k(F)=k0 exp[FΔx/(kBT)]，无力时k0=0.005000s\textasciicircum{}-1。施加单分子张力100.0pN、反应势垒位移Δx=0.2000nm，温度298.15K，Boltzmann常数1.380649×10\textasciicircum{}-23J/K。求外力降低势垒的J/分子、无量纲势垒变化、加速倍数、受力后k和半衰期，指出pN·nm与SI J的换算。\par
\medskip\noindent\textbf{分步解答}
\begin{enumerate}
\item 将100pN化为10\textasciicircum{}-10N，0.2nm化为2×10\textasciicircum{}-10m。
\item 力做功FΔx=2e-20J/分子。
\item 无量纲W/(kBT)=4.858608429。
\item 加速倍数exp[W/(kBT)]=128.8447807。
\item 受力后k=k0×倍数=0.6442239034s\textasciicircum{}-1。
\item 对应半衰期t1/2=ln2/k=1.075941419s；如果Δx随力变化，这个简化指数律不再通用。
\item 独立核验全部反应物账本、状态概率和动力学物理解残差：0.00e+0、0.00e+0、0.00e+0。
\end{enumerate}
\noindent\textbf{答案摘要}\par
答案1 = 2e-20；答案2 = 4.85860843；答案3 = 128.844781；答案4 = 0.644223903；答案5 = 1.07594142\par
\noindent\textbf{模型前提：}
仅假设题给真实物料与动力学机制；跨辑去重、实验真实性及人类专家核定尚待完成。\par
\noindent\textcolor{red!65!black}{[跨辑去重未完成 · 草稿]} \par
\medskip\hrule\medskip
\Needspace{7\baselineskip}
\subsubsection*{UC-ADV9-O89 · 酸碱双催化有机水解U形pH—速率曲线的最低点}
\noindent\textbf{题目}\par
有机酯恒温水解按kobs=kH[H+]+kOH[OH−]+k0进行，kH=1.000×10\textasciicircum{}4L mol\textasciicircum{}-1s\textasciicircum{}-1、kOH=1.000×10\textasciicircum{}5L mol\textasciicircum{}-1s\textasciicircum{}-1、k0=1.000×10\textasciicircum{}-5s\textasciicircum{}-1、Kw=1.000×10\textasciicircum{}-14，强缓冲固定pH且无其它催化剂。求kobs在全部正h范围最小时的h及pH、此时酸碱两条反应的各自速率贡献、最低总kobs与半衰期，证明求得的是全局极小而非极大。\par
\medskip\noindent\textbf{分步解答}
\begin{enumerate}
\item 令h=[H+]，Kw/h=[OH−]，得kobs(h)=kHh+kOHKw/h+k0。
\item 对h求导k′=kH−kOHKw/h²=0，得到h*=√(kOHKw/kH)=3.16227766e-7M。
\item 最优pH=−log10(h*)=6.5。
\item 此时酸与碱两项均为0.00316227766s\textasciicircum{}-1。
\item 最低kobs=0.00633455532s\textasciicircum{}-1，半衰期=109.423179s。
\item 因为二阶导k''=2kOHKw/h³>0，该点是唯一全局最小值。
\item 独立核验全部反应物账本、状态概率和动力学物理解残差：-4.34e-19、0.00e+0、0.00e+0。
\end{enumerate}
\noindent\textbf{答案摘要}\par
答案1 = 3.16227766e-07；答案2 = 6.5；答案3 = 0.00633455532；答案4 = 109.423179\par
\noindent\textbf{模型前提：}
仅假设题给真实物料与动力学机制；跨辑去重、实验真实性及人类专家核定尚待完成。\par
\noindent\textcolor{red!65!black}{[跨辑去重未完成 · 草稿]} \par
\medskip\hrule\medskip
\Needspace{7\baselineskip}
\subsubsection*{UC-ADV9-O90 · Woodward–Hoffmann电环化热/光条件下旋转方式的四格判别}
\noindent\textbf{题目}\par
按基态轨道对称守恒规则，具有4π电子的开链共轭体系热电环化允许conrotatory（同向旋转），具有6π电子的体系热电环化允许disrotatory（反向）；单光子激发时选择规则反转。某题要求对4π与6π体系分别在热与光下判断是否conrotatory，输出四个二元指示量（允许conrotatory为1，否则0），求四格中允许conrotatory的总数和光化学相对热化学的翻转次数，并说明允许不意味着实际产率为100\%。\par
\medskip\noindent\textbf{分步解答}
\begin{enumerate}
\item 4π热电环化由基态前线轨道对称性要求同向旋转，conrotatory指标1。
\item 6π热电环化要求反向旋转，conrotatory指标0。
\item 4π光激发轨道占据翻转选择规则，conrotatory指标0。
\item 6π光化学则允许同向旋转，conrotatory指标1。
\item 四项中恰两项允许conrotatory=2，与热条件相比光照翻转两项。
\item 轨道对称性规则只说明对称允许/禁阻，不自动给出反应速度、产物光稳态或精确立体收率。
\item 核对速率、统计权重、电子数和热力学方程的独立关系，残差：0.00e+0、0.00e+0、0.00e+0。
\end{enumerate}
\noindent\textbf{答案摘要}\par
答案1 = 1；答案2 = 0；答案3 = 0；答案4 = 1；答案5 = 2\par
\noindent\textbf{模型前提：}
题给理想化机理及条件；不等同真实实验或专家科学与原创审定。\par
\noindent\textcolor{red!65!black}{[跨辑去重未完成 · 草稿]} \par
\medskip\hrule\medskip
\Needspace{7\baselineskip}
\subsubsection*{UC-ADV9-O91 · 正交保护基的选择性脱保护与竞争误反应存活率}
\noindent\textbf{题目}\par
假想双保护官能团化合物同时带酸敏A保护基和基本不受酸影响的B保护基；酸处理下A去保护一级kA=0.120min\textasciicircum{}-1、B意外脱保护一级kB=0.002000min\textasciicircum{}-1，两事件独立。处理15.00min后目标化合物要求A已经脱掉、B仍保留。求A脱保护分数、B存活分数、目标双条件分数、B错误脱保护分数，并求若要至少90\% A去保护所需最短反应时间是否还能保证B存活率≥95\%。\par
\medskip\noindent\textbf{分步解答}
\begin{enumerate}
\item A去保护概率PA=1−exp(−kA t)=0.834701112。
\item B仍保留概率PB=exp(−kB t)=0.970445534。
\item 由于独立，目标A去且B留比例PA PB=0.810031966。
\item 错误脱掉B的边际概率1−PB=0.0295544665。
\item A去保护率至少90\%需要t≥ln(10)/kA=19.1882091min。
\item 在这个最短时刻B存活率exp[−kB ln(10)/kA]=0.962350626；与0.95比较判断是否同时达标。
\item 核对速率、统计权重、电子数和热力学方程的独立关系，残差：0.00e+0、0.00e+0、0.00e+0。
\end{enumerate}
\noindent\textbf{答案摘要}\par
答案1 = 0.834701112；答案2 = 0.970445534；答案3 = 0.810031966；答案4 = 19.1882091；答案5 = 0.962350626\par
\noindent\textbf{模型前提：}
题给理想化机理及条件；不等同真实实验或专家科学与原创审定。\par
\noindent\textcolor{red!65!black}{[跨辑去重未完成 · 草稿]} \par
\medskip\hrule\medskip
\Needspace{7\baselineskip}
\subsubsection*{UC-ADV9-O92 · 协同二位点有机酶Hill反应速率和半饱和底物浓度}
\noindent\textbf{题目}\par
某有机催化酶的表观活化满足Hill式v=Vmax[S]\textasciicircum{}n/(Khalf\textasciicircum{}n+[S]\textasciicircum{}n)，给定协同指数n=2、Vmax=0.01000M/s、Khalf=0.02000M。自由底物浓度由连续进料库固定为0.03000M，求占据分数θ、反应初速、达到0.9Vmax所需底物浓度和在S=Khalf时的对数斜率dln[v/(Vmax−v)]/dlnS，指出n不是必然实际结合位点数。\par
\medskip\noindent\textbf{分步解答}
\begin{enumerate}
\item Hill占位θ=S²/(K²+S²)=0.692307692。
\item 速率v=Vmaxθ=0.00692307692M/s。
\item θ=0.9时S²/(K²+S²)=0.9，解S=3K=0.06M。
\item Hill作图log[θ/(1−θ)]=nlog(S/K)，其对数斜率=n=2。
\item S=K处θ=0.5、速率Vmax/2=0.005M/s。
\item 经验Hill系数只是协同性指标，不能直接宣称恰好有两个结构上独立的结合位点。
\item 核对速率、统计权重、电子数和热力学方程的独立关系，残差：0.00e+0、0.00e+0、-1.11e-16。
\end{enumerate}
\noindent\textbf{答案摘要}\par
答案1 = 0.692307692；答案2 = 0.00692307692；答案3 = 0.06；答案4 = 2；答案5 = 0.005\par
\noindent\textbf{模型前提：}
题给理想化机理及条件；不等同真实实验或专家科学与原创审定。\par
\noindent\textcolor{red!65!black}{[跨辑去重未完成 · 草稿]} \par
\medskip\hrule\medskip
\Needspace{7\baselineskip}
\subsubsection*{UC-ADV9-O93 · Ping–Pong Bi–Bi双底物催化与单底物饱和的相互作用}
\noindent\textbf{题目}\par
假想有机转氨酶遵循Ping–Pong Bi–Bi初速模型v=Vmax[A][B]/(KB[A]+KA[B]+[A][B])，Vmax=0.01200M/s，KA=0.02000M、KB=0.01500M，两个自由底物A=0.03000M、B=0.02500M由外部保持恒定。求初速、B翻倍后的初速、速率提高倍率及当B→∞时A固定条件下的极限初速，解释这与简单两底物速率相乘有何不同。\par
\medskip\noindent\textbf{分步解答}
\begin{enumerate}
\item 写分母D=KB A+KA B+A B=0.0017M²。
\item 代入给定浓度，初速v=Vmax A B/D=0.00529411765M/s。
\item B翻倍至0.05000M，按同一分母结构得v2=0.00610169492M/s。
\item 提高倍率v2/v=1.15254237。
\item B无限大时分子分母同除B，极限v∞=Vmax A/(KA+A)=0.0072M/s。
\item Ping–Pong速率式假设特定共价/酶态循环且取初速，不能套作所有双底物酶催化机理。
\item 核对速率、统计权重、电子数和热力学方程的独立关系，残差：0.00e+0、0.00e+0、0.00e+0。
\end{enumerate}
\noindent\textbf{答案摘要}\par
答案1 = 0.00529411765；答案2 = 0.00610169492；答案3 = 1.15254237；答案4 = 0.0072\par
\noindent\textbf{模型前提：}
题给理想化机理及条件；不等同真实实验或专家科学与原创审定。\par
\noindent\textcolor{red!65!black}{[跨辑去重未完成 · 草稿]} \par
\medskip\hrule\medskip
\Needspace{7\baselineskip}
\subsubsection*{UC-ADV9-O94 · RAFT/ATRP受控聚合活性链分数与单体增长速率}
\noindent\textbf{题目}\par
假想受控自由基聚合链端存在休眠D和活性R两态，激活一级速率ka=0.003000s\textasciicircum{}-1、失活一级kd=0.02700s\textasciicircum{}-1。链端总浓度CP=0.002000M，单体由大库保持[M]=0.5000M，增长常数kp=150L mol\textasciicircum{}-1s\textasciicircum{}-1，仅活性链可传播。求活性链分数、活性自由基浓度、瞬时单体消耗速率及100s若浓度全程恒定每条初始链平均增加的重复单元数；说明这一近似不含终止与单体耗竭。\par
\medskip\noindent\textbf{分步解答}
\begin{enumerate}
\item 快链端平衡活性分数pR=ka/(ka+kd)=0.1。
\item 活性链浓度[R]=CP pR=0.0002M。
\item 增长速率v=kp[M][R]=0.015M/s。
\item 100s单体消耗浓度Δ[M]=v×100=1.5M，由大库补充才可保持给定浓度。
\item 每mol原始链增加的平均重复单元数=ΔM/CP=750。
\item 若活性链持续与休眠态交换，链增长数是时间平均；终止和浓度变化要求更复杂的动态方程。
\item 核对速率、统计权重、电子数和热力学方程的独立关系，残差：0.00e+0、0.00e+0、0.00e+0。
\end{enumerate}
\noindent\textbf{答案摘要}\par
答案1 = 0.1；答案2 = 0.0002；答案3 = 0.015；答案4 = 750\par
\noindent\textbf{模型前提：}
题给理想化机理及条件；不等同真实实验或专家科学与原创审定。\par
\noindent\textcolor{red!65!black}{[跨辑去重未完成 · 草稿]} \par
\medskip\hrule\medskip
\Needspace{7\baselineskip}
\subsubsection*{UC-ADV9-O95 · 环闭复分解平衡与乙烯分压作为外部状态变量}
\noindent\textbf{题目}\par
假想有机二烯A在金属卡宾催化剂下发生环闭复分解A⇌C+E，其中E为乙烯气体。温度体积固定，外部排气与乙烯分压控制使乙烯的无量纲活度aE=0.01000恒定，溶液中的A、C活度按c/1M近似。条件平衡常数K=[C]aE/[A]=20.00，初始[A]=0.03000M、[C]=0，求平衡环状产物C、残余A、转化率和若乙烯活度提高十倍的转化率；解释必须把气相乙烯的活动度与液相物料分开。\par
\medskip\noindent\textbf{分步解答}
\begin{enumerate}
\item 设x=[C]，液相底物A=A0−x，气体乙烯活度aE由外部库钳定。
\item 平衡关系K=x aE/(A0−x)，解x=KA0/(K+aE)=0.0299850075M。
\item 剩余[A]=A0−x=0.0000149925037M。
\item 转化率x/A0=K/(K+aE)=0.99950025，低乙烯活度推动环闭。
\item 将乙烯活度提高至0.1000，转化率降为K/(K+0.1000)=0.995024876。
\item 若乙烯气体封闭积累，其活度不再恒定，须同时联立气相状态方程及乙烯生成量。
\item 核对速率、统计权重、电子数和热力学方程的独立关系，残差：-3.30e-12、0.00e+0、0.00e+0。
\end{enumerate}
\noindent\textbf{答案摘要}\par
答案1 = 0.0299850075；答案2 = 1.49925037e-05；答案3 = 0.99950025；答案4 = 0.995024876\par
\noindent\textbf{模型前提：}
题给理想化机理及条件；不等同真实实验或专家科学与原创审定。\par
\noindent\textcolor{red!65!black}{[跨辑去重未完成 · 草稿]} \par
\medskip\hrule\medskip
\Needspace{7\baselineskip}
\subsubsection*{UC-ADV9-O96 · 双峰分子量分布聚合物的Mn、Mw与分散度不同权重}
\noindent\textbf{题目}\par
一个假想有机聚合物样品由两种近似单分散链组成，按链条数量计80.00\%的链摩尔质量M1=10000Da、20.00\%的链M2=50000Da。求数均摩尔质量Mn、重均Mw、分散度Đ=Mw/Mn以及M2分子在样品总质量中的贡献比例，证明不能把Mw直接做数分数加权的算术平均。\par
\medskip\noindent\textbf{分步解答}
\begin{enumerate}
\item 按数量分数，Mn=pM1+(1−p)M2=18000Da。
\item 重均定义Mw=Σn\_iM\_i²/Σn\_iM\_i，故Mw=32222.2222Da。
\item 分散度Đ=Mw/Mn=1.79012346。
\item 大分子M2总质量比例w2=(1−p)M2/Mn=0.555555556。
\item 尽管M2只占20\%链条数量，却贡献更高质量分数，因此显著拉高Mw。
\item 如果样品还有宽连续分布，则应对整个分子量分布的零、一、二阶矩积分。
\item 核对速率、统计权重、电子数和热力学方程的独立关系，残差：1.82e-12、0.00e+0、-3.64e-12。
\end{enumerate}
\noindent\textbf{答案摘要}\par
答案1 = 18000；答案2 = 32222.2222；答案3 = 1.79012346；答案4 = 0.555555556\par
\noindent\textbf{模型前提：}
题给理想化机理及条件；不等同真实实验或专家科学与原创审定。\par
\noindent\textcolor{red!65!black}{[跨辑去重未完成 · 草稿]} \par
\medskip\hrule\medskip
\Needspace{7\baselineskip}
\subsubsection*{UC-ADV9-O97 · 有机电合成法拉第效率与化学选择性的两级收率}
\noindent\textbf{题目}\par
假想有机电还原电解池恒流0.5000A、运行600.0s，目标还原每mol底物转移2mol电子。扣除副电极反应后可参与有机底物还原的法拉第效率ηF=0.8000，其中真正生成目标产品的化学选择性S=0.7500，剩余有机还原副产物也按2电子计。求总电荷、有效有机反应电子mol、目标产品mol和有机副产物mol，以及目标产品相对总电荷的总体电子利用率，注意电流效率和化学选择性不能混同。\par
\medskip\noindent\textbf{分步解答}
\begin{enumerate}
\item 总电荷Q=It=0.5000×600=300C。
\item 电子总mol ne=Q/F=0.0031092809。
\item 有机底物反应得到的电子mol=ηF ne=0.00248742472。
\item 目标产品电子份额再乘选择性S，转2电子得到目标产物mol=0.000932784269。
\item 有机副产物mol=0.00031092809，按题设亦消耗2电子。
\item 总体目标电子利用率ηtarget=ηF S=0.6，两级效率必须相乘。
\item 核对速率、统计权重、电子数和热力学方程的独立关系，残差：0.00e+0、0.00e+0、0.00e+0。
\end{enumerate}
\noindent\textbf{答案摘要}\par
答案1 = 300；答案2 = 0.00248742472；答案3 = 0.000932784269；答案4 = 0.00031092809；答案5 = 0.6\par
\noindent\textbf{模型前提：}
题给理想化机理及条件；不等同真实实验或专家科学与原创审定。\par
\noindent\textcolor{red!65!black}{[跨辑去重未完成 · 草稿]} \par
\medskip\hrule\medskip
\subsection*{物理化学}
\addcontentsline{toc}{subsection}{第9辑 / 物理化学}
\Needspace{7\baselineskip}
\subsubsection*{UC-ADV9-P73 · van der Waals气体的单相压缩因子与临界参数交叉复核}
\noindent\textbf{题目}\par
假想单组分气体服从van der Waals状态方程P=RT/(Vm−b)−a/Vm²，其中a=0.3600Pa·m6/mol²、b=4.300×10\textasciicircum{}-5m³/mol，T=350.0K、P=5.000MPa。求可行的Vm、压缩因子Z=PVm/(RT)、理论临界温度Tc=8a/(27Rb)、临界压强Pc=a/(27b²)，检查Vm>b以及给定T相对临界温度的位置，说明不能套理想气体Vm=RT/P替代真实根。\par
\medskip\noindent\textbf{分步解答}
\begin{enumerate}
\item 先计算临界温度Tc=8a/(27Rb)=298.350028K。
\item T=350K高于模型临界温度，因此按题定单相气体分支寻找Vm>b的物理解。
\item 在Vm∈(b,0.020)m³/mol对Pcalc(Vm)−5MPa二分求根，得Vm=0.0004912852661m³/mol。
\item 压缩因子Z=PVm/(RT)=0.8441147994。
\item 临界压强Pc=a/(27b²)=7.211105102MPa。
\item 因Z≠1，有限排除体积和分子吸引不可由理想气体单项反算。
\item 独立检验能量、离子活动度或非理想输运的正向约束，残差：3.73e-16、0.00e+0、0.00e+0。
\end{enumerate}
\noindent\textbf{答案摘要}\par
答案1 = 0.000491285266；答案2 = 0.844114799；答案3 = 298.350028；答案4 = 7.2111051\par
\noindent\textbf{模型前提：}
题设热力学、理想或截断近似有效；实际实验及人工科学/新颖性审核尚未完成。\par
\noindent\textcolor{red!65!black}{[跨辑去重未完成 · 草稿]} \par
\medskip\hrule\medskip
\Needspace{7\baselineskip}
\subsubsection*{UC-ADV9-P74 · 混合二价—一价电解质的离子强度和Davies活度校正}
\noindent\textbf{题目}\par
298.15K理想充分解离的混合盐溶液含MgCl2=0.01000M、NaCl=0.02000M，不发生配位或沉淀。按I=1/2Σci zi²计算离子强度，并以Davies近似log10γi=−0.509zi²[√I/(1+√I)−0.3I]计算γMg2+、γCl−和纯MgCl2组分的平均离子活度系数γ±=(γMg γCl²)\textasciicircum{}(1/3)。求各值并解释不能用只含MgCl2的浓度算混合液的I。\par
\medskip\noindent\textbf{分步解答}
\begin{enumerate}
\item 完全解离给Mg2+=0.01000M、Na+=0.02000M、Cl−=0.04000M。
\item 离子强度I=1/2(4×.010+.020+.040)=0.05M。
\item Davies共同括号Q=√I/(1+√I)−0.3I=0.1677439976。
\item 二价Mg活度系数γMg=0.4554840995，单价氯γCl=0.8215203109。
\item MgCl2平均系数γ±=(γMg γCl²)\textasciicircum{}(1/3)=0.6748956212。
\item Davies是稀溶液近似，若I更高应检验适用范围，不能称其全浓度精确。
\item 独立检验能量、离子活动度或非理想输运的正向约束，残差：0.00e+0、0.00e+0、0.00e+0。
\end{enumerate}
\noindent\textbf{答案摘要}\par
答案1 = 0.05；答案2 = 0.4554841；答案3 = 0.821520311；答案4 = 0.674895621\par
\noindent\textbf{模型前提：}
题设热力学、理想或截断近似有效；实际实验及人工科学/新颖性审核尚未完成。\par
\noindent\textcolor{red!65!black}{[跨辑去重未完成 · 草稿]} \par
\medskip\hrule\medskip
\Needspace{7\baselineskip}
\subsubsection*{UC-ADV9-P75 · Kohlrausch独立离子迁移定律与离子迁移数的交叉预测}
\noindent\textbf{题目}\par
无限稀释水溶液的限制摩尔电导给定Λ∞(NaCl)=126.5、Λ∞(KCl)=149.9、Λ∞(KNO3)=144.9，单位均S·cm²/mol。已知Na+离子限制电导λNa=50.10同单位，阴离子Cl−由NaCl的差值获得。用Kohlrausch离子独立迁移定律求Λ∞(NaNO3)、NaCl中Na+迁移数和通过该电解质截面的1.000A电流对应Na+电荷迁移mol/s，并解释迁移数只表示载流贡献而非反应消耗速率。\par
\medskip\noindent\textbf{分步解答}
\begin{enumerate}
\item 无限稀释的盐电导等于各离子限制电导之和。
\item 组合相消得到Λ∞NaNO3=ΛNaCl+ΛKNO3−ΛKCl=121.5S·cm²/mol。
\item NaCl中Cl−离子电导λCl=126.5−50.1=76.4。
\item Na+迁移数tNa=λNa/(λNa+λCl)=0.3960474308。
\item 1.000A的电荷流速是I/F mol电子当量/s，由Na+贡献的电荷当量流=0.00000410474237mol/s。
\item 迁移数随浓度、溶剂和温度改变，有限浓度不能无条件沿用无限稀释值。
\item 独立检验能量、离子活动度或非理想输运的正向约束，残差：0.00e+0、0.00e+0、5.55e-17。
\end{enumerate}
\noindent\textbf{答案摘要}\par
答案1 = 121.5；答案2 = 0.396047431；答案3 = 4.10474237e-06；答案4 = 76.4\par
\noindent\textbf{模型前提：}
题设热力学、理想或截断近似有效；实际实验及人工科学/新颖性审核尚未完成。\par
\noindent\textcolor{red!65!black}{[跨辑去重未完成 · 草稿]} \par
\medskip\hrule\medskip
\Needspace{7\baselineskip}
\subsubsection*{UC-ADV9-P76 · Goldman恒场膜电势的双一价阳离子渗透率竞争}
\noindent\textbf{题目}\par
理想选择性膜只允许K+与Na+扩散，在298.15K遵循无净电流Goldman简化式Vm=(RT/F)ln[(PK[K]out+PNa[Na]out)/(PK[K]in+PNa[Na]in)]。膜内[K]=140mM、[Na]=10mM，膜外[K]=5mM、[Na]=145mM，渗透率比PK:PNa=20:1。求Vm、只允许K+时的Nernst电势EK、两者偏移以及加权浓度比，解释膜对K更通透不意味着Vm必严格等于EK。\par
\medskip\noindent\textbf{分步解答}
\begin{enumerate}
\item 恒场式先求外侧加权项20×5+145=245mM比例单位。
\item 内侧加权项20×140+10=2810mM。
\item 加权浓度比Rconc=0.0871886121，对数是无量纲。
\item 膜电势Vm=(RT/F)ln(Rconc)=-0.0626817113V。
\item 纯钾离子Nernst电势EK=(RT/F)ln(5/140)=-0.08561292803V。
\item 偏移Vm−EK=0.02293121673V；简单式依赖独立离子渗透和恒电场假设，不能无限扩展到多价离子。
\item 独立检验能量、离子活动度或非理想输运的正向约束，残差：0.00e+0、0.00e+0、0.00e+0。
\end{enumerate}
\noindent\textbf{答案摘要}\par
答案1 = 0.0871886121；答案2 = -0.0626817113；答案3 = -0.085612928；答案4 = 0.0229312167\par
\noindent\textbf{模型前提：}
题设热力学、理想或截断近似有效；实际实验及人工科学/新颖性审核尚未完成。\par
\noindent\textcolor{red!65!black}{[跨辑去重未完成 · 草稿]} \par
\medskip\hrule\medskip
\Needspace{7\baselineskip}
\subsubsection*{UC-ADV9-P77 · 曲面液滴Kelvin蒸气压与Laplace压力的联立推导}
\noindent\textbf{题目}\par
298.15K某纯液体表面张力γ=0.07200N/m、摩尔体积Vm=18.00cm³/mol。形成半径r=10.00nm的凸球形液滴，假设γ和Vm不随曲率变化，理想蒸气满足Kelvin式ln(p(r)/p∞)=2γVm/(rRT)。求Laplace内外压差、Kelvin指数、饱和蒸气压比和把曲面误当凹液面的预测方向，解释该模型在纳米尺度需考虑表面张力尺寸依赖。\par
\medskip\noindent\textbf{分步解答}
\begin{enumerate}
\item 球形液滴压差ΔP=2γ/r=14.4MPa。
\item 摩尔体积18.00cm³/mol=18×10\textasciicircum{}-6m³/mol。
\item Kelvin无量纲指数2γVm/(rRT)=0.1045601021。
\item 凸液滴蒸气压比p(r)/p∞=exp(exponent)=1.110222118。
\item 凸曲率提高平衡蒸气压；对于相同曲率大小的凹液面，指数符号反向。
\item 半径若接近分子尺寸，连续表面张力与理想蒸气假设可能失败。
\item 独立检验能量、离子活动度或非理想输运的正向约束，残差：0.00e+0、-5.55e-17、0.00e+0。
\end{enumerate}
\noindent\textbf{答案摘要}\par
答案1 = 14.4；答案2 = 0.104560102；答案3 = 1.11022212；答案4 = 0.900720661\par
\noindent\textbf{模型前提：}
题设热力学、理想或截断近似有效；实际实验及人工科学/新颖性审核尚未完成。\par
\noindent\textcolor{red!65!black}{[跨辑去重未完成 · 草稿]} \par
\medskip\hrule\medskip
\Needspace{7\baselineskip}
\subsubsection*{UC-ADV9-P78 · 非离子表面活性剂的Gibbs吸附量与分子截面积}
\noindent\textbf{题目}\par
298.15K非离子表面活性剂水溶液，在活度a1=0.01000时平衡表面张力γ1=0.07200N/m，a2=0.1000时γ2=0.06800N/m，活度无量纲。假设该范围γ对ln a近似线性，Gibbs吸附式Γ=−(1/RT)dγ/dln a。求表面过量Γ(mol/m²)、每平方米分子数、单分子平均占据面积(nm²)，并解释这里面积推断隐含单层紧密占据假设。\par
\medskip\noindent\textbf{分步解答}
\begin{enumerate}
\item 取两点近似dγ/dln a=(γ2−γ1)/ln(a2/a1)=-0.001737177928N/m。
\item Gibbs式Γ=−斜率/(RT)=7.007696813e-7mol/m²。
\item 表面分子数密度N=Γ NA=422013366100000000m\textasciicircum{}-2。
\item 每个分子平均面积A=1/N=2.369593194nm²。
\item γ随活度增加而下降使Γ为正，说明溶质在界面富集。
\item 真正的单分子面积还取决于界面覆盖结构与聚集，Gibbs表面过量不总等于完整单层位点数量。
\item 独立检验能量、离子活动度或非理想输运的正向约束，残差：-8.67e-19、0.00e+0、0.00e+0。
\end{enumerate}
\noindent\textbf{答案摘要}\par
答案1 = 7.00769681e-07；答案2 = 4.22013366e+17；答案3 = 2.36959319\par
\noindent\textbf{模型前提：}
题设热力学、理想或截断近似有效；实际实验及人工科学/新颖性审核尚未完成。\par
\noindent\textcolor{red!65!black}{[跨辑去重未完成 · 草稿]} \par
\medskip\hrule\medskip
\Needspace{7\baselineskip}
\subsubsection*{UC-ADV9-P79 · BET线性化与单层吸附量表面积的反演}
\noindent\textbf{题目}\par
一种多孔材料满足BET线性式y=p/[v(p0−p)]=1/(vm C)+(C−1)/(vm C)×(p/p0)。给定独立拟合线截距b=1/(20.00×80.00)、斜率m=79/(20.00×80.00)，v单位cm³(STP)/g，标准态1mol气体占22414cm³，N2分子占据面积0.1620nm²。求单层容量vm、BET参数C、在p/p0=0.2000时吸附体积和BET比表面积，并明确多层模型区别于Langmuir单层。\par
\medskip\noindent\textbf{分步解答}
\begin{enumerate}
\item 由BET线性式b=1/(vm C)、m=(C−1)/(vm C)。
\item 因此vm=1/(b+m)=20cm³(STP)/g。
\item C=1+m/b=80。
\item 取x=p/p0=.2，y=b+mx=0.0105，由y=x/[v(1−x)]反解v=23.80952381cm³/g。
\item 比表面积S=vm/22414×NA×0.162nm²，换为m²/g得到87.05155734m²/g。
\item BET只在合适相对压力区间拟合并假设多层吸附，微孔填充可能使计算表面积偏离真实。
\item 独立检验能量、离子活动度或非理想输运的正向约束，残差：0.00e+0、0.00e+0、0.00e+0。
\end{enumerate}
\noindent\textbf{答案摘要}\par
答案1 = 20；答案2 = 80；答案3 = 23.8095238；答案4 = 87.0515573\par
\noindent\textbf{模型前提：}
题设热力学、理想或截断近似有效；实际实验及人工科学/新颖性审核尚未完成。\par
\noindent\textcolor{red!65!black}{[跨辑去重未完成 · 草稿]} \par
\medskip\hrule\medskip
\Needspace{7\baselineskip}
\subsubsection*{UC-ADV9-P80 · Debye介电弛豫复介电常数与损耗峰频率}
\noindent\textbf{题目}\par
某假想极性液体在GHz范围遵循单Debye介电弛豫ε*(ω)=ε∞+(εs−ε∞)/(1+iωτ)，静态εs=80.00、高频ε∞=4.000、弛豫时间τ=8.000ps。求损耗峰角频率和普通频率、该频率下ε'、ε''（采用ε*=ε'−iε''）及损耗角正切tanδ=ε''/ε'，并解释频率太低或太高时损耗项为何趋零。\par
\medskip\noindent\textbf{分步解答}
\begin{enumerate}
\item 对于单Debye过程，损耗ε''=Δε ωτ/[1+(ωτ)²]。
\item 其峰值发生ωτ=1，所以ωpeak=1/τ=125000000000rad/s。
\item 普通频率fpeak=ω/(2π)=19.89436789GHz。
\item 在ωτ=1处实部ε'=ε∞+Δε/2=42。
\item 损耗ε''=Δε/2=38，损耗角正切=0.9047619048。
\item 当ωτ≪1或≫1，ε''分别随ωτ和1/(ωτ)衰减，真实液体多个弛豫时间会展宽峰形。
\item 独立检验能量、离子活动度或非理想输运的正向约束，残差：0.00e+0、0.00e+0、0.00e+0。
\end{enumerate}
\noindent\textbf{答案摘要}\par
答案1 = 19.8943679；答案2 = 42；答案3 = 38；答案4 = 0.904761905\par
\noindent\textbf{模型前提：}
题设热力学、理想或截断近似有效；实际实验及人工科学/新颖性审核尚未完成。\par
\noindent\textcolor{red!65!black}{[跨辑去重未完成 · 草稿]} \par
\medskip\hrule\medskip
\end{document}
